\documentclass[../main.tex]{subfiles}
\begin{document}
%==========================================TIÊU ĐỀ TẬP========================================
\begin{table}[h]
    \begin{adjustwidth}{-1cm}{}
        \begin{tabular}{>{\centering\arraybackslash}p{6cm}|>{\raggedright\arraybackslash}p{14cm}}
            \multirow{5}{6cm}
            % Cột bên trái: ÉPISODE và số tập
            {\\[-40pt]
                \flaregothic\fontsize{20pt}{20pt}\selectfont \centering
                \color{eptype}HISTOIRE\color{black}\\
                \dawnland\fontsize{72pt}{72pt}\selectfont
                \color{epnum}4
                \\[18pt]
                % \flaregothic\fontsize{14pt}{14pt}\selectfont
                % \color{epattr}BIENVENUE, \uline{\color{Banpire} Mel} !
                % \flaregothic\fontsize{18pt}{18pt}\selectfont
                % \color{epattr}ÉPISODE PERDU
                \color{black}
            }
             &
            % Dòng thứ nhất: tên tập tiếng Nhật
            {\fotskip\fontsize{18pt}{18pt}\selectfont \ruby{異}{い}\ruby{世}{せ}\ruby{界}{かい}\ruby{火}{ひ}\ruby{川}{かわ}\ruby{姉}{し}\ruby{妹}{まい}} \\[6pt]
            % Dòng thứ hai: tên tập tiếng Anh
             & {\cafeteria\fontsize{18pt}{18pt}\selectfont Sisters from Another World}                                       \\[4pt]
            % Dòng thứ ba: tên tập tiếng Việt
             & {\vnmsans\fontsize{18pt}{18pt}\selectfont Hai chị em từ thế giới khác}                                                   \\[8pt]
            % Dòng thứ tư: Nhân vật xuất hiện
             & {\caviar{\fontsize{14pt}{14pt}\selectfont Nhân vật: \emp{kuon1} \emp{kaigou1} \emp{suzuki1}}}                               \\[8pt]
        \end{tabular}
    \end{adjustwidth}
\end{table}
%===========================================NỘI DUNG==========================================
\color{black}

\pov{Hikawa Kuon}

Chúng tôi tiếp tục bước đi. Tôi, cô gái có vẻ ngoài lạnh lùng mà bên trong ẩn giấu bao bão giông mang tên ``chị gái,'' và một cô bé phàm ăn nhút nhát nhưng lại có trái tim ấm áp nhất trong cái nơi kỳ dị này.

Suzuki-chan vẫn bám lấy lưng tôi như một chú gấu koala nhỏ, đôi khi ngọ nguậy vì nhột hoặc do mỏi chân (mặc dù cô bé chỉ thấp hơn tôi cỡ một gang tay), còn Kaigou thì bước đi bên cạnh, tay khoanh lại, gương mặt chẳng rõ đang khó chịu hay đang cố giả vờ nghiêm nghị.

Tôi nói vài câu bâng quơ để phá vỡ bầu không khí: ``\vie{Này, Suzuki-chan, em ăn nhiều thế, có bao giờ bị đau dạ dày chưa?}''

``\vie{Ơ\dots\ chưa ạ}'' cô bé thỏ thẻ, xong khúc khích. ``\vie{Dạ dày em khỏe mà!}''

``\vie{Tốt rồi, không thì giờ anh sẽ bị kiện vì cõng người bệnh mất.}''

Cô bé cười to, còn cô chị chỉ hừ một tiếng rõ to. Cô quay sang lườm tôi: ``\vie{Cõng con bé thì cõng, còn nhiều lời quá.}''

Tôi nhún vai, liếc sang cô: ``\vie{Tôi đâu có bắt cô đi bộ bằng mồm đâu?}''

``\vie{Cái cậu này\dots!}''

Lần này thì tôi nghe thấy tiếng ``khẹc'' rõ ràng từ phía bên kia. Đúng là chị em ruột có khác, nói chuyện với cả hai người này mà tôi cảm giác như bị bao vây giữa hai cơn bão; một cơn bão xấu hổ nhưng nhí nhảnh và một cơn bão bực bội đang học cách làm dịu bản thân lại.

Lúc mới bước chân vào ``The Void'', tôi cứ tưởng mình sẽ phát điên, tất cả chỉ bởi cái tin nhắn chết tiệt từ giáo viên bộ môn. Những thứ vô lý, lơ lửng, trôi dạt như một giấc mơ chưa tỉnh cứ thế ập vào đầu tôi. Không điểm bắt đầu, không mục tiêu rõ ràng, chỉ có những câu hỏi vô tận trong đầu và hàng tá vật thể kỳ quặc không ngừng hiện ra.

Tưởng rầng tôi sẽ phải đi bộ một cách vô định ở một nơi rỗng tuếch này \emph{một mình} với cơn tức giận vẫn còn âm ỉ trong lòng, thì giờ đây, tôi có hai người đồng hành.

Một người khiến tôi không biết nên thở dài hay khâm phục vì lòng can đảm của cô ta, người kia thì nhẹ nhàng như một mảnh khăn gió thoảng qua, nhưng lại đủ mạnh để gắn kết cả ba lại với nhau.

Tôi không ngờ rằng mình lại đi xa đến mức này chỉ vì\dots\ một sự sắp đặt lơ đễnh và đầy ngẫu hứng của giáo viên bộ môn chúng tôi.

Nhưng nhờ thế, tôi gặp được Kaigou. Và cả Suzuki.

Những câu chuyện của họ, những phản ứng, ánh mắt, những lần chạm nhẹ như vô tình nhưng đầy cảm xúc\dots\ tất cả đều khiến tôi thấy một điều: ở cái nơi kỳ quái và bất định này, sự sống không nằm ở không gian, mà nằm trong con người.

Và có lẽ giờ đây, với tâm trạng có phần thoải mái hon của hai người con gái tôi có thể tìm hiểu kỹ hơn về con người, về xuất thân, tất cả mọi thứ về họ.

Biết đâu đấy, tôi với Kaigou có đặc điểm gì thú vị thì sao nhỉ?

\ct

Sau một hồi chúng tôi đi bộ mệt mỏi, cả ba chúng tôi quyết định ngồi nghỉ chân tại một cái lán (thực chất chỉ là mấy mỏm đá được dựng lên để lấy làm ghế ngồi). Cả ba chúng tôi sau đó kể những câu chuyện đời thường mà ba chúng tôi hay gặp phải. Có chuyện thì khiến tôi phì cười vì độ ngớ ngẩn, có chuyện thì tôi ngẩn tò te vì không hiểu họ đang nói tới cái gì.

Tôi không nhớ chính xác ai là người mở lời trước. Có thể là tôi, có thể là Suzuki-chan, mà cũng có thể là Kaigou, người đang lườm tôi vì cái gì đó, và tôi nhân cơ hội đấy để đổi chủ đề.

``\vie{Này, tôi hỏi thật nhé\dots\ Thế giới của hai người từng sống như thế nào vậy?}''

Kaigou liếc tôi, hơi nhíu mày như thể định phản ứng phòng thủ, nhưng lại thôi. Có lẽ vì cô cũng hiểu tôi không hỏi để thẩm vấn, mà là vì thực lòng tò mò.

Suzuki-chan lúc này đã ngồi xuống bên cạnh tôi, vẫn còn hơi đỏ mặt vì chuyện vừa nãy, nhưng mắt thì long lanh như thể sẵn sàng kể chuyện cổ tích. Cô bé nhìn sang chị mình như đợi tín hiệu.

Cô chị khẽ gật và bắt đầu câu chuyện. ``\vie{Chúng tôi đến từ một nơi, mà theo cách gọi của em tôi thì là một ``phiên bản khác'' của thế giới mà cậu đã biết. Về cơ bản, vẫn là Etsunan. Vẫn là điện thoại thông minh, vẫn là giao thông nườm nượp, vẫn là quán ăn gia đình mở 24/7. Nhưng\dots}''

Cô ngừng lại một chút, đánh mắt về cô em gái. Cô bé tiếp lời với vẻ thận trọng: ``\vie{Nhưng có một điều khác biệt lớn, Kuon-nii-san ạ. Ở thế giới của bọn em, tỷ lệ nam nữ là khoảng\dots\ ba-trên-một (\(3 : 1\)).}''

Tôi nhíu mày. Ba-trên-một?

``\vie{Ý em là\dots\ ba người đàn ông thì mới có một người phụ nữ?}'' tôi hỏi lại, hơi kinh ngạc.

Kaigou lắc đầu. ``\vie{Ngược lại. Ba phụ nữ, một đàn ông.}

``\vie{À, ra thế.}'' Tôi gãi đầu. Mọi thứ bắt đầu rõ hơn rồi. Rất nhiều thứ.

Kaigou thở dài, mắt không nhìn tôi mà hướng về một nơi vô định phía trước, rồi thở dài tiếp tục: ``\vie{Vì vậy, đàn ông ở nơi đó được chiều chuộng, đề cao, và dễ dàng nắm quyền hơn. Không phải tất cả, nhưng đủ để tạo nên một xã hội bất cân bằng. Nhất là với những kẻ biết lợi dụng.}''

Suzuki-chan nắm lấy chị mình, khẽ gật đầu đồng tình: ``\vie{Onee-sama từng gặp phải nhiều chuyện không vui với họ, nên mới cảnh giác như thế ạ.}''

Tôi nhìn sang Kaigou. Tôi không cần nghe cô ấy kể thêm cũng biết chuyện đã xảy ra nghiêm trọng đến mức nào.

Cô chị lớn quay lại nhìn tôi, vẻ mặt bình thản nhưng có gì đó trong mắt như một lớp băng chưa tan hết. ``\vie{Ở nơi đó, con gái như em tôi phải học cách trốn, né tránh, hoặc cam chịu. Tôi thì không chấp nhận điều đó. Tôi bảo vệ nó, khỏi bất kỳ thằng nào mon men lại gần, bằng mọi giá.}

Tôi im lặng một lúc, rồi lên tiếng: ``\vie{Kể cả là với người tốt?}''

``\vie{\dots Tôi không có thời gian để phân biệt ai tốt ai xấu. Khi lớn lên giữa bầy thú hoang, sẽ không ai tin con nào là `thú cưng thuần hóa' cả,}'' cô nàng thở dài mạnh hơn. ``\vie{Vậy nên cách dễ nhất là cứ gom chúng về một mối thôi.}''

Câu nói ấy đập thẳng vào lòng tôi, nhưng tôi không thấy giận. Chỉ thấy một cái gì đó\dots\ buồn. Vì một người từng phải sống như vậy. Sống trong xã hội nơi mà người tốt hiếm như vàng như thế, đến tôi còn thấy bất ổn, huống gì là họ.

Suzuki-chan khẽ lên tiếng: ``\vie{Nhưng bọn em vẫn có một ngôi nhà nhỏ. Một căn phòng bé tẹo ở khu chung cư cũ. Em rất thích cái đệm nằm cạnh cửa sổ\dots\ mỗi tối hai chị em hay nằm đó, nói mấy chuyện ngớ ngẩn. Bọn em cũng từng có tivi, nhưng lâu lắm rồi không xem.}''

Giờ thì tôi hiểu vì sao họ biết mọi thứ về đồ dùng của tôi, vì nơi họ sống giống y hệt thế giới của tôi. Nhưng tôi cũng hiểu tại sao cô chị thì chẳng bao giờ tin tôi, còn cô em thì ngập ngừng như đang sống với một vết sẹo dài trên tâm hồn.

Hai chị em nhà này mang theo một phần thế giới vỡ nát của họ bước vào nơi không có thực này, cùng với vết thương không dễ gì lành lại.

Nhưng có lẽ, chỉ cần đi cùng nhau, và cùng tôi, vết thương ấy sẽ ít đau hơn một chút.

Tôi chợt hỏi nhỏ, xoáy sâu hơn chút về gia thế hai người: ``\vie{Thế còn bố mẹ hai người thì sao?}''

Bỗng nhiên, cả hai chị em cùng im lặng. Người chị lên tiếng, giọng trầm hơn: ``\vie{Chúng tôi đã sống một mình từ năm Suzuki lên tám. Còn tôi thì lúc đó mới mười tám.}''

``\vie{Chuyện gì xảy ra sao?}'' Tôi hỏi, gần như theo phản xạ. Câu nói tưởng chừng vô hại, chỉ đơn giản là tiếp nối cuộc trò chuyện, nhưng tôi nhận ra mình vừa lỡ lời ngay khi thấy ánh mắt hai người họ khựng lại.

Kaigou quay mặt đi, còn Suzuki-chan thì cụp mắt xuống, hai tay siết lấy vạt áo. Dường như tôi đã vô tình chạm vào một vết thương chưa bao giờ thực sự lành.

Im lặng một lúc, cuối cùng Kaigou mới lên tiếng. Giọng cô đều đều, nhưng chất chứa những thứ còn nặng hơn cả nỗi buồn. Một thứ gì đó mà chưa hoàn toàn nhạt phai.

``\vie{Mẹ bọn tôi\dots\ mất trong một tai nạn giao thông. Lúc đó bà đang chở hàng bằng xe máy cho một cửa hàng bán buôn. Trời đổ mưa, đường trơn, xe tải trượt bánh và rồi\dots}

Cô ngừng lại, không cần nói hết câu. Tôi cũng không cần nghe nốt phần còn lại để hiểu.

Suzuki-chan khẽ lẩm bẩm: ``\vie{Em còn nhớ lúc đó, cảnh sát đến tận nhà. Em mới năm tuổi. Chị thì đi học, không có ở đó. Em đứng trong hành lang, không hiểu gì cả. Họ chỉ hỏi: `Có người thân lớn tuổi không?', rồi gọi chị về.}''

``\vie{\dots}'' Tôi không cần phải lên tiếng, nhưng cũng đủ để mừng tượng được chuyện đó tệ hại ra sao.

Tôi nhìn sang cô chị, và thấy trong mắt cô một thứ gì đó nghẹn lại. Không phải nước mắt, một cái gì nén lại từ rất lâu.

``\vie{Ba năm sau, đến lượt bố chúng tôi,} cô nói tiếp. ``\vie{Tai nạn lao động. Ông ấy là thợ cơ khí, làm trong nhà máy chế tạo. Có một lần, máy ép bị trục trặc, bố tôi ngó vào kiểm tra, và thế là xong.}

Tôi điếng người một lúc, rồi nuốt nước bọt, cổ họng khô khốc.

``\vie{Tôi xin lỗi\dots\ Tôi không nên hỏi điều đó\dots} tôi khẽ nói, cúi đầu. ``\vie{Tôi không biết là hai chị em cô đã phải chịu nhiều đau đớn đến vậy\dots}

Kaigou không đáp. Cô chỉ khoanh tay, mắt nhìn về phía xa, vẻ như đang giữ cho mình không gục xuống vì một chuyện đã trôi qua bảy năm.

Tôi thở dài, thầm ước rằng không động tới câu hỏi này.

Suzuki-chan mím môi, nhưng rồi quay sang tôi, gượng cười, cố giấu đi nỗi đau mà tôi vô tình khơi ra: ``\vie{Không sao đâu Kuon-nii-san ạ. Nếu không kể, chắc hai chị em sẽ cứ giữ mãi trong lòng như một cái hộp khóa chặt, không có một ai bên cạnh để tâm sự đâu.}''

Tôi không biết nói gì hơn. Chỉ khẽ đặt tay lên vai cô bé, ấn nhẹ một cái. ``\vie{Hai người đã mạnh mẽ hơn anh tưởng rất nhiều rồi đấy.}''

Câu nói đó làm Kaigou hơi liếc nhìn tôi, lần đầu tiên không mang ý phòng thủ hay gay gắt. Cô không đáp, chỉ lặng lẽ kéo nhẹ cô em gái lại gần mình như muốn san sẻ nỗi buồn.

Suzuki-chan tựa vào vai chị, khẽ nói, nhắc lại những kỷ niệm đẹp khi gia đình họ vẫn còn đông đủ để gạt đi nỗi buồn ấy. ``\vie{Mẹ hay kể chuyện cho em nghe mỗi tối. Hồi đó còn có bố hay vào bếp trổ tài nấu ăn, tuy dở tệ, nhưng cả nhà ăn vẫn vui ạ.}''

``\vie{Ông ấy chẳng bao giờ cho đủ nước vào nồi cơm điện,}'' cô chị bật ra một nụ cười hiếm hoi. ``\vie{Có lần còn quên gạt nút nấu, đến lúc mở ra vẫn thấy gạo sống nguyên.}''

Cả ba chúng tôi cùng cười nhẹ. Một nụ cười ngắn, nhưng đủ để xoa dịu những gì vừa nặng nề kéo đến.

Tôi nhìn hai chị em họ, một người vẫn còn mang nét trẻ con ngây thơ, một người đã phải mang gánh nặng gia đình trên vai khi chưa tròn hai mươi. Tôi không dám nói rằng mình hiểu hoàn toàn họ, rằng tôi muốn hỏi ``làm sao có thể''?

Nhưng ít ra, tôi có thể lắng nghe.

Tôi im lặng một lúc lâu, để họ có thời gian với chính ký ức của mình. Trong những lúc như thế này, sự im lặng cũng là một cách thể hiện sự tôn trọng.

``\vie{Sau đó thì sao?}'' tôi hỏi tiếp, một cách chậm rãi. ``\vie{Hai người đã sống ra sao khi không còn bố mẹ?}''

Kaigou thở ra môt hơi nặng nề, tay kéo nhẹ chiếc áo khoác mà tôi cho cô ấy mượn. Một câu hỏi không mấy dễ chịu, nhưng dường như cô ấy không từ chối trả lời.

``\vie{Tôi phải nghỉ học một thời gian để đi làm thêm. Mấy công việc lặt vặt thôi: phục vụ quán ăn, phát tờ rơi, làm đêm ở cửa hàng tiện lợi\dots\ Cứ có tiền là làm. Lúc đó Suzu vẫn còn nhỏ, tôi không muốn em phải chuyển vào trại trẻ mồ côi.}''

Suzuki-chan cúi đầu, tay mân mê gấu váy, giọng lí nhí tiếp lời từ chị: ``\vie{Em biết chị đã vất vả nhiều lắm. Nhưng em không biết phải làm gì để giúp cả\dots}''

Người chị xoa đầu cô em gái, giọng hiền hòa: ``\vie{Chị chỉ cần Suzu học hành đàng hoàng là đủ rồi.}''

Tôi nhìn hai người, lòng chợt nhói lên một chút. Có những thứ mà tôi tưởng là hiển nhiên, như một bữa cơm đủ đầy với cả gia đình, hay chỉ đơn giản là một ngày Chủ Nhật nằm nghỉ ngơi ở nhà, thì với họ đó lại là những thứ xa xỉ, thậm chí có thể gọi là chỉ có thể nhìn thấy trong mơ.

``\vie{Vậy tiền có đủ sống không?}'' tôi hỏi, không giấu được sự quan tâm.

Dĩ nhiên, câu trả lời nhận được từ Kaigou không khiến tôi ngạc nhiên: một chữ ``Không'' đầy dứt khoát. ``\vie{Chúng tôi sống nhờ trợ cấp xã hội một phần, nhưng vẫn thiếu trước hụt sau. Có lần tôi bị cắt lương vì khách than phiền, lần khác thì cửa hàng đóng cửa đột xuất. Suzuki phải ăn mì gói liền cả tuần. Tôi từng phải gõ cửa xin cơm nguội từ nhà bên cạnh\dots}''

Tôi siết chặt nắm tay, cố gắng không để cảm xúc tràn ra ngoài mặt. Ở cái tuổi mà người ta còn chưa biết viết đơn xin việc, thậm chí nếu không vì cái đơn ngẫu nhiên nào đó trên mạng, chưa chắc tôi đã có được một công việc ổn định, thì Kaigou đã gồng cả một sự trách nhiệm lớn trên vai.

Y hệt như những gì mà bà nội tôi hay kể hồi bà ấy còn minh mẫn, khi kịch bản cũng tương tự như vậy: một người mẹ phải nuôi năm người con, trong một gia cảnh nghèo khó, khi cái ăn còn không được đảm bảo.

Nhưng đó là chuyện của thời bao cấp ngày xưa, còn thời hiện đại bây giờ thì khác. Dẫu cho là thời kỳ nào, họ cũng là những người đáng được tôn trọng.

``\vie{Thế còn chuyện học?}'' tôi bèn quay sang hỏi Suzuki-chan. ``\vie{Sao em lại học lớp Tám?}''

Suzuki đỏ mặt một chút, không vì xấu hổ mà vì hơi ngại: ``Em bị nghỉ học mấy lần vì gia đình không đủ điều kiện đóng học phí. Có lúc em chuyển trường, có lúc chị bảo em ở nhà chờ ổn định đã\dots\ rồi dần dần bị chậm hai năm so với các bạn cùng tuổi ạ.''

Tôi lại nhớ đến khoảng thời gian học cấp hai, những lần mà tôi lúc nào cũng đón tiền học muộn nhất cả lớp, tới mức mà giáo viên chủ nhiệm của tôi phải đóng hộ. Nếu không vì tôi học tốt và trở thành một trong những học sinh ưu tú của lớp hồi đó, có lẽ tôi cũng sẽ rơi vào tình cảnh như vậy.

Tôi gật đầu, không nói gì ngay. Một phần trong tôi thấy giận, không phải giận họ, mà giận cái hệ thống đã không thể giúp họ đúng lúc, đã đẩy họ vào cuộc sống đầy rẫy sự khó khăn đó.

Nhưng phần lớn còn lại là cảm giác\dots\ tự thấy bản thân nhỏ bé. Chỉ mới hơn hai năm trước thôi, tôi mới chỉ là một kẻ thích nằm ngủ ở nhà, ngồi chơi điện tử cả ngày rồi ăn. Hầu như chưa bao giờ phải động tay động chân làm gì để kiếm tiền phụ giúp cả.

Một người từng nghĩ mình cũng trải qua không ít khó khăn, lại chưa từng phải lo ngày mai có gì ăn, hay có còn chỗ ngủ không.

Ngẫm nghĩ lại, tôi cảm thấy mình còn thua xa cả một đứa trẻ học lớp Tám, chứ chưa nói là một cô chị bằng tuổi tôi.

``\vie{Vậy mà em vẫn giữ được nụ cười đó, Suzuki-chan?}'' tôi hỏi, hơi nghiêng đầu.

Cô bé cười khẽ, nụ cười nhẹ tênh nhưng không yếu đuối: ``\vie{Em nghĩ\dots\ nếu em cũng gục xuống và bỏ cuộc như thế, thì Onee-sama biết bám vào đâu? Cả hai chị em chỉ còn có nhau để sống nương tựa nhau thôi chứ.}''

Tôi nhìn sang Kaigou. Cô hơi ngoảnh mặt đi, nhưng tôi bắt được ánh nhìn lặng lẽ mang chút gì ấm áp trong mắt cô.

Không biết từ khi nào, tôi không còn nhìn họ như những người xa lạ nữa.

Tôi bắt đâu cảm thấy\dots\ đây là một gia đình.

Dù không hoàn hảo, dù đã vỡ vụn, nhưng vẫn cố gắng từng ngày để nhặt nhạnh từng mảnh vụn còn sót lại, để tiếp tục sống, tiếp tục tồn tại ở một xã hội đầy rẫy sự nghi kỵ và sự hỗn loạn từ những kẻ cơ hội, những kẻ dâm đãng ở đó.

``\vie{Hai đứa từng sống ở đâu vậy?}'' tôi tiếp tục hỏi, ngồi bó gói trên mỏm đá. ``\vie{Chính xác hơn là sau khi trở thành trẻ mồ côi ấy.}''

Cô chị trả lời ngay, giọng khô như khói sương: ``\vie{Một khu nhà trọ tự gọi là `chung cư mini.' Phòng nhỏ, trần thấp, tường mốc, không cách âm, và phải dùng nhà vệ sinh chung.}''

``\vie{Chung cư mini?}'' tôi nghiêng đầu. Ý của họ là mấy cái nhà trọ cao tầng cho những người có lao động thu nhập thấp thuê à? Tôi đã từng nghe những chuyện không hay liên quan tới mấy cái kiểu nhà như thế, nên tôi cũng mừng tượng được cuộc sống của hai cô gái đó như thế nào.

Suzuki-chan nở một nụ cười trầm buồn, rồi tiếp lời cô chị: ``\vie{Chúng em sống tầng ba, không có thang máy. Cầu thang lúc nào cũng bẩn, nhớp nhớp nước, rác bịt kín dưới từng bậc. Mỗi lần mưa là nước tràn vào nhà. Nấm mốc mọc cả trên trần.}''

Kaigou gật đầu, tiếp tục kể khổ: ``\vie{Lúc thuê, chủ trọ bảo là có wifi, có nước nóng, có bảo vệ. Nhưng thực tế thì chẳng có cái nào ra hồn cả. Wifi yếu đến mức không tải nổi email, nước nóng chỉ hoạt động vài giờ mỗi ngày, còn bảo vệ thì chỉ xuất hiện đầu mỗi tháng lúc thu tiền nhà thôi.}''

Lúc nghe hai chị em kể, tôi tưởng tượng tới viễn cảnh cứ mỗi mùa mưa bão, Suzuki-chan phải ôm cặp lội qua hành lang tối om, mùi ẩm mốc vấn vương vào tóc, bước qua từng vũng nước loang lổ không được dọn sạch trong đôi giày rách.

``\vie{Nhiều lúc bọn em còn bị ghẻ do lội nước nhiều quá đó, Kuon-nii-san ạ. Bọn em phải tự bôi thuốc cho nhau để kìm cự nữa cơ.}''

Tôi buông một cái thở dài. Cuộc sống của hai chị em họ đúng là khốn khó thật. Một góc nhỏ nào đó trong tim tôi như chùng xuống.

Tôi tiếp tục hỏi tới những thứ khó nói hơn: ``\vie{Có chuyện gì bất công mà hai người nhớ nhất không? Kiểu, cái gì khiến hai người phẫn uất nhất?}''

Cả hai chị em im lặng trong một lúc. Người chị lớn đáp, giọng vẫn trầm và có chút kiềm nén: ``\vie{Một hôm tôi vừa tan ca đêm, định tranh thủ tạt qua siêu thị mua đồ ăn sáng cho Suzu, thì có một gã trong khu trọ - tôi vẫn nhớ mặt hắn ta - tự tiện bước vào phòng, vì cửa hỏng khóa.}''

Tôi bỗng thấy lạnh sống lưng. Chuyện một kẻ trộm đột nhập vào chung cư để trộm cắp vặt, thậm chí là cướp tài sản giá trị tôi đã từng nghe nhiều, nhưng đột nhập vào một căn phòng hầu như không có gì\dots\ lại là một chuyện khác.

Suzuki-chan siết chặt bàn tay, khẽ run lên bần bật vì sợ: ``\vie{Lúc đó em đang ngủ một mình ạ. Em chỉ biết chuyện này qua lời kể của Onee-sama thôi.}''

``\vie{Hắn không làm gì em chứ?}'' tôi bật ra, nét mặt hiện rõ vẻ lo lắng.

Cô bé lắc đầu: ``\vie{Em không sao. May mà chị về kip.}''

``\vie{Cậu biết đấy, sau đó tôi\dots}'' cô chị gái nắm tay thành nắm đấm. ``\vie{\dots đã đánh hắn. Không nặng tay lắm, nhưng đủ để hắn sợ mà tránh xa bọn tôi.}''

Tôi trầm ngâm. Không phải vì tôi không biết những chuyện như vậy xảy ra, mà là vì chúng lại xảy ra với chính họ, hai con người vừa yếu ớt vừa kiên cường. ``\vie{Vậy cô đã giải quyết chuyện này cho quản lý như thế nào?}''

``\vie{Tôi đã báo chuyện này lên quản lý rồi, nhưng họ bảo không có camera, không có chứng cứ, nên họ không giải quyết. Họ còn bảo rằng `ở đâu mà chẳng có mấy chuyện này,' rồi khuyên tôi `đóng cửa cho kỹ' thay vì làm lớn chuyện. Họ cũng đã hứa là sẽ sửa khóa nữa, nhưng họ có bao giờ làm đâu.}''

Tôi không nói được gì. Tôi chỉ biết rút ra hộp sữa từ ba-lô, đưa cho Suzuki-chan. Cô bé nhận lấy, mỉm cười gượng gạo.

``\vie{Từ đó tôi cắt luôn giao tiếp với hàng xóm,} Kaigou chốt lại, mắt lạnh như thép. ``\vie{Niềm tin một khi đập vỡ thì không thể hàn gắn được dễ dàng.}''

``\textit{Và thế là cô vô tình làm ngơ luôn những người tốt tính thực sự còn gì,}'' tôi thầm nghĩ. Nhưng tôi hiểu tại sao cô ấy phải làm vậy, ít nhất là theo bối cảnh xã hội và cách mà thế giới của hai chị em hoạt động như thế nào. Tôi không thể nào áp quy chuẩn nơi tôi ở với nơi họ ở được.

Trong lòng tôi, cái tên Kaigou, người từng chỉ là một cô gái khó chịu vừa hung dữ vừa hay nghi ngờ, bắt đầu được khắc lại dưới một góc nhìn khác. Một người chị mạnh mẽ, quyết liệt, sẵn sàng bảo vệ em gái mình khỏi cả thế giới nếu cần.

Mà nếu ngẫm nghĩ lại\dots\ thì việc cô ấy muốn đấm tôi khi lần đầu gặp lại là điều dễ hiểu.

Những câu chuyện mà Kaigou và Suzuki-chan kể vẫn còn văng vẳng trong đầu tôi: nặng trĩu và cay đắng.

Thế nhưng, dường như chính họ lại là những người gỡ nút thắt đó trước tiên.

Cô em gái ngẩng đầu lên, ánh mắt hướng vào khoảng tối trước mặt. Giọng cô bé vang nhẹ: ``\vie{Nhưng mà\dots\ chỗ đó cũng không phải lúc nào cũng tệ đâu.}''

Tôi nhìn sang. Cô bé đang nắm tay chị gái mình, và nở một nụ cười nhỏ.

``\vie{Hồi đó, mỗi tối, bọn em thường chia nhau một cái quạt máy bé tí, chạy còn không nổi mức hai. Nhưng nếu nghiêng nó đúng hướng\dots}'' cô bé xòe tay ra diễn tả, ``\vie{\dots thì cả hai chị em sẽ được mát một lượt. Bọn em cứ xoay xoay cho đến khi đúng `góc vàng', rồi nằm cười vì thấy mình khôn ra.}''

``\vie{Đúng là khi gặp khó thì chị em mình luôn có cách nhỉ,}'' Kaigou bật cười nhẹ, lần đầu tiên từ khi tôi biết cô ấy. Một nụ cười rất chân thật, rất người.

``\vie{Còn nhớ hôm mưa to không?}'' cô liếc mắt xuống. ``\vie{Hôm em lấy thau ra hứng nước dột, rồi bày thêm mấy cốc giấy làm thành trò `thuỷ chiến' nho nhỏ. Em bày trò cho nó vui, nhưng rồi lại bị sặc vì ngã vào nước\dots}''

``\vie{\dots chị thì cứ cười ngặt nghẽo, còn em thì ướt như chuột lột!}'' Suzuki-chan phụ họa, cười rúc rích như thể chuyện đó mới chỉ xảy ra ngày hôm qua vậy.

Tôi cũng bật cười. Không lớn, nhưng là nụ cười thật lòng đầu tiên kể từ khi tôi tới nơi quái quỷ này. Lồng ngực tôi bỗng thấy nhẹ đi, như thể áp lực bám víu suốt từ lúc lạc vào một nơi rỗng không vừa tan đi một mảng lớn.

``\vie{Anh dám cá mấy trò ngốc nghếch ngớ ngẩn kia là niềm vui nhỏ nhoi của hai người nhỉ?}''

Cô chị gật đầu, ánh mắt dịu hơn hẳn: ``\vie{Nếu không tự tạo niềm vui, thì chẳng ai làm cho mình đâu. Những lúc được cười thật lòng, dù chỉ một chút thôi, là những thứ giúp chúng tôi lạc quan hơn.}''

Cô em gái mỉm cười, ngước nhìn tôi: ``\vie{Có mấy hôm trời lạnh, bọn em gom mấy cái khăn cũ lại làm tổ trên giường, bật điện thoại nghe truyện ma rẻ tiền trên mạng. Rồi giật bắn người lên vì\dots\ tiếng hàng xóm tắm phòng bên. Hài kinh khủng!}'' và bật cười.

Tôi phá ra cười lần nữa. Không phải vì những gì họ trải qua, mà vì cách họ kể có pha một sự tếu táo và mộc mạc. Hai chị em ấy tuy trải qua nhiều thứ chẳng ai mong, nhưng vẫn giữ được một thứ gì đó trong sáng và tử tế.

Có thể chính điều đó mới là thứ mạnh mẽ nhất giúp họ tồn tại được ở thế giới khắc nghiệt đó.

``\vie{Nhưng mà, tôi thắc mắc một điều\dots}'' tôi cất tiếng khi cuộc trò chuyện lắng xuống. ``\vie{Hai người không có họ hàng nào sao? Bạn bè, người thân đâu mà hai người phải khổ sở như vậy?}''

Suzuki lặng đi một thoáng. Còn Kaigou thì nhìn thẳng về phía trước, ánh mắt chùng xuống. ``\vie{Có chứ. Đúng hơn, là đã từng. Hồi tôi còn nhỏ, nhà Hikawa có hẳn một dòng họ lớn, đủ thứ chú bác, dì cô, họ hàng xa gần. Nhưng\dots}

Cô ngập ngừng, như đang cố gắng lục lại ký ức đã bụi mờ từ lâu. ``\vie{Hồi tôi lên chín tuổi, gia đình bên nội xảy ra một vụ tranh chấp tài sản.}''

Tôi khẽ nhíu mày. Đoạn này tôi đã phần nào đoán trước được, thậm chí gia đình tôi còn đang manh nha một vụ như thế. Nhưng nghe từ chính miệng cô ấy, vẫn khiến tôi chững lại, càng không khiến tôi rùng mình.

``\vie{Bắt đầu từ chuyện thừa kế đất đai của ông bà cố. Rồi kéo theo hàng loạt mâu thuẫn, nghi kỵ, kiện tụng\dots\ Sau đó, cả nhà tan vỡ, mỗi người đi một nẻo,}'' cô chị tiếp tục, không hề tỏ vẻ nao núng gì, như thể việc đó không liên quan gì tới cô.

Kaigou cười xòa, tiếp tục: ``\vie{Cậu biết rồi đấy, mỗi người về một ngả, có người dọn đi tỉnh khác, có người thì cạch mặt nhau, tuyệt giao luôn. Họ hàng vốn đã không thân, giờ thì chẳng còn gì để gọi là ruột thịt nữa.}''

\vie{Lúc mẹ chúng em còn sống, bà vẫn giữ liên lạc với một vài người. Nhưng sau khi bố mẹ qua đời rồi\dots}'' Suzuki-chan nói nhỏ, ánh mắt vẻ đượm buồn, ``\vie{\dots bọn em bị bỏ mặc hoàn toàn.}''

``\vie{Tôi từng đi tìm vài người cô chú mà tôi tiếp xúc thường xuyên, nhưng\dots}'' Kaigou lắc đầu, thở ra một hơi dài đầy chán nản, ``\vie{chẳng ai muốn vướng thêm rắc rối, nhất là khi tôi mang theo một đứa em gái nhỏ. Với họ, tôi chỉ là cái tên gợi lại những cuộc cãi vã trong quá khứ mà họ muốn quên.}''

Tôi im lặng. Môi khẽ mím lại. Những tưởng sẽ nghe một câu đại loại như ``mất liên lạc'' hay ``bận công việc'' đơn thuần, ai ngờ là cả một câu chuyện đổ vỡ như vậy. Dựa theo câu chuyện mà hai cô nàng kể, tôi dám chắc mọi chuyện không đơn giản nếu chỉ xét vẻ bề ngoài.

Nhưng tôi quyết định không xoáy sâu về vấn đề đó. Vấn đề gia đình luôn là thứ mà tôi không nhúng tay vào, đặc biệt là chuyện nội bộ của nhà khác.

``\vie{Cả đàn ông lẫn phụ nữ mà tôi gặp khi đi cầu viện đều khiến tôi mệt mỏi. Người thì cười khẩy, kẻ thì nhìn tôi như món nợ thừa thãi từ đời nào. Có lần, tôi còn bị một ông chú đề nghị đưa Suzu đi đổi lấy tiền trợ giúp.}''

Tôi ôm đầu, lắc đầu ngán ngẩm thay cho Kaigou. Tôi nghe rõ từng từ, mà vẫn không tin nổi mình đang nghe những chuyện này từ miệng một cô gái chỉ mới hai mươi lăm tuổi.

``\vie{Từ đó trở đi, tôi không còn tin ai nữa. Tôi chỉ biết, nếu không giữ lấy em mình, thì sẽ không ai giữ hộ cả.}''

Tôi nhìn sang Suzuki-chan. Cô bé không khóc, nhưng ánh mắt ngân ngấn nước. Nỗi đau đã được giấu kín từ lâu, giờ được thốt ra tới một người mà họ đã bắt đầu mở lòng tin tưởng.

Giờ tôi đã hiểu. Hiểu vì sao cô ấy luôn cố mạnh mẽ đến vô lý. Vì sao cô bé hiền lành kia lại không thể dựa vào ai ngoài chị mình. Và vì sao cái họ Hikawa\dots\ dường như chỉ còn lại hai người duy nhất trong cái thế giới khắc nghiệt ấy.

\ct

Sau một lúc nghỉ ngơi ngắn và ngẫm nghĩ lại mọi chuyện, ba người chúng tôi tiếp tục cuộc hành trình không có hồi kết trong khoảng không bí ẩn.

Tôi im lặng bước bên hai chị em họ, mắt dán về phía trước, nhưng đầu óc thì đang xoay vòng, cố lắp ghép lại những thông tin mà hai chị em kia vừa thổ lộ.

Câu chuyện về thế giới của họ đang dần xếp thành từng mảnh ghép trong đầu tôi. Một xã hội nơi tỷ lệ nam nữ mất cân bằng nghiêm trọng - theo nghĩa ngược lại mà chúng ta thường biết - nơi mà phụ nữ bị biến thành trò mua vui và bị khinh thường tới rẻ mạt. Một gia đình tan vỡ, không phải vì lỗi của bất kỳ ai trong số họ, mà là do tai nạn, nghịch cảnh, rồi tranh chấp gia tộc. Một cô chị mới mười tám tuổi đã gánh cả cuộc đời lên vai. Một cô em học muộn hai năm, rèn luyện tính tiết kiệm triệt để tới độ nhịn tiểu cả ngày trời.

Tôi nhớ lại ánh mắt của Kaigou lúc kể về gia đình mình. Giọng cô ấy khi nhắc tới người bố, một giọng bình tĩnh đến mức tôi biết chắc nó đã được rèn giũa qua hàng trăm đêm nén tiếng khóc vì mất đi người thân yêu.

Tôi không phải loại người dễ xúc động khi nghe những câu chuyện như thế, nhưng không thể giả vờ như mình không biết gì nữa.

Suốt quãng thời gian trong The Void, tôi từng nghĩ đây chỉ là một tai nạn, một vụ dịch chuyển không gian đơn thuần. Nhưng sau từng câu từng chữ họ nói, tôi nhận ra: nếu có một nơi nào đó mà người ta gọi là ``không còn gì để mất nữa'', thì chính là nơi họ đã bước ra.

Tôi thở nhẹ.

Nếu là tôi, rơi vào hoàn cảnh ấy, tôi sẽ làm gì?

Không. Câu hỏi không nên như thế.

Nếu họ quay về được thế giới của mình, họ sẽ làm gì tiếp theo?

Tôi liếc sang Kaigou. Cô ấy đang lặng lẽ bước đi, một tay giữ chặt lấy túi đeo hông, một tay cầm lấy bàn tay nhỏ của Suzuki-chan.

Cô em gái thì mỉm cười, bám víu lấy lưng tôi như khỉ, ngân nga giai điệu nào đó mà có lẽ tôi đã quên.

Tôi vẫn chưa thể đưa ra quyết định.

Không phải vì tôi không biết nên làm gì. Tôi thực sự vẫn chưa rõ, nếu quay trở lại cái nơi được gọi là ``nhà'' kia, họ sẽ có tương lai nào để bám víu không?

Tôi cần biết điều đó, trước khi đưa ra lời đề nghị cuối cùng.

Tôi vừa định mở miệng thì Kaigou đột ngột dừng bước.

``\vie{Cậu định hỏi, `nếu hai chị em tôi quay trở về, chúng tôi sẽ làm gì', đúng không?}''

Tôi khựng lại, đôi chân gần như đông cứng giữa nhịp bước.

\dots Sao cô ấy biết?

Tôi nhìn vào mắt cô nàng. Nó không toát ra vẻ lạnh lẽo, không cảnh giác cao độ, mà là một thứ gì đó thấu hiểu đến mức gần như vô thức.

``\vie{Chắc là tôi đoán được,}'' cô cười nhẹ, nhún vai, rồi quay đầu nhìn sang phía em gái mình. ``\vie{Tôi từng đoán đúng suy nghĩ của Suzu suốt bao năm. Có vẻ như\dots\ tôi cũng làm thế được với cậu. Một chút thôi.}''

Tôi bật ra một tiếng ``hừm'' trong cổ họng. Không rõ là cười hay là thở dài. ``\vie{Nhạy bén quá cũng chẳng phải điều tốt đâu.}''

``\vie{Tôi biết. Nhưng nếu tôi không nhạy bén, chúng tôi đã chết thẳng cẳng từ thuở nào rồi.}''

Giọng nói ấy lần này, không giấu được sự trầm xuống. Tôi cảm tưởng một câu chuyện buồn nữa sắp sửa được tiết lộ.

Tôi im lặng, như để dành cho cô thời gian.

``\vie{\dots Cái đệm mà Suzu nhắc đến ấy, chiếc mà cậu rơi trúng khi vào đây.}''

Tôi gật đầu. Tôi vẫn nhớ cảm giác mùi khét lẹt và vết cháy đen trên mép vải.

Kaigou hít vào một hơi, rồi tiếp tục giãi bày, nhìn thẳng vào mắt tôi. ``\vie{Thứ đó chính là thứ cuối cùng còn lại trong căn phòng trước khi mọi thứ bị nuốt vào cái hố đó. Lửa cháy chưa được bao lâu thì trời\dots\ không biết gọi là gì nữa. Nó mở ra. Không gian xé rách. Như một cái miệng không đáy. Và chúng tôi bị kéo vào đây.}''

Cô dừng lại ở đó.

Như vậy, Kaigou cùng Suzuki-chan đã nhìn thấy thứ đó. Vòng xoáy bí ẩn kia.

Tại sao nó lại xuất hiện tại nơi đó? Phải, tôi biết là cái vòng xoáy kia xuất hiện bởi thứ gọi là ``lỗ hổng ma thuật'' mà Shion-chan đã từng nhắc tới ba ngày trước.

Nhưng ở nơi của hai chị em kia? Tôi không tài nào biết được.

Tôi hiểu ngay đây không chỉ là một câu chuyện đơn thuần. Đây là ký ức. Một đoạn ký ức nặng nề đến mức chỉ cần chạm vào là cũng khiến cô nghẹt thở.

``\vie{Kể lại không dễ đâu. Nhưng tôi sẽ nói. Có lẽ\dots\ cậu cũng cần biết một chút.}''

Tôi gật nhẹ, không nói gì thêm.

Có lẽ, chính tôi cũng đang cần nghe điều đó. Không chỉ để hiểu họ, mà còn để có thể hiểu chính cái nơi điên rồ này, ``The Void''.

\pov{Hikawa Kaigou}

\begin{center}
    \uline{Hồi tưởng\\10:30 tối - khoảng ba ngày trước.}
\end{center}

Tôi nghĩ\dots\ nó xảy ra vào khoảng ba đêm trước. Chắc vậy.

\dots Không, tôi không chắc nữa. Trong cái nơi như thế này, nơi mà không có ngày đêm, không có thời gian. Đồng hồ trên máy của tôi dừng lại ở lúc 0 giờ 11 phút sáng. Còn khi mọi thứ bắt đầu\dots\ là đâu đó hơn 10 giờ rưỡi tối.

Chúng tôi - tôi và Suzu - vừa ăn tối xong.

\il

Tối hôm ấy không có gì đặc biệt. Tôi làm ca tối ở một tiệm tiện lợi ở loanh quanh \ruby{Seishun}{Thanh Xuân}. Cũng không quá xa đâu, chỉ cách khoảng mười phút đi bộ về nhà thôi.

Lúc tôi về tới nhà thì đồng hồ đã chỉ gần mười rưỡi, khoảng thời gian mà bình thường Suzu đã chìm vào giấc ngủ, nhưng vào những ngày hè, con bé sẵn sàng thức để chờ tôi trở về.

Vào tới căn phòng với cái lưng mỏi nhừ, mắt nặng trĩu vì buồn ngủ, khi bước chân vào nhà, tôi vẫn thấy con bé lăng xăng lau dọn sàn, trên trán buộc một chiếc khăn như mọi lần. Ánh đèn vàng vọt của khu trọ chiếu lên tấm lưng bé nhỏ ấy khiến tôi vừa thấy ấm, vừa thấy chạnh lòng.

``\vie{Suzu vẫn chưa ngủ à?}'' tôi khẽ trố mắt ngạc nhiên hỏi.

``\vie{Đợi chị về ăn cơm đó!}'' cô bé quay lại, nở một nụ cười toe toét, tay vẫn đang giũ giẻ lau.

Tôi ngồi xuống, cởi giày, ngồi bệt xuống đệm, vừa xoa cổ chân tê rần, vừa rên rỉ: ``\vie{Mệt quớ đi\~{}}'' Tôi ngước nhìn về phía bàn ăn đã được soạn sẵn từ khi nào, trong đó có một âu cơm đã đặt ở đó, trông còn khá mới.

``\vie{Có cơm à?}'' tôi hỏi như bản năng.

``\vie{Ưm! Cơm mà cô hàng xóm phòng bên cạnh cho đấy ạ! Em mang đồ xuống rửa thì cô bảo dư cơm, cho thêm phần. Có hộp cá nướng chị thích nữa này!}''

Tôi khựng lại một nhịp.

Tôi không phải kiểu người tin người dễ dàng, càng không thích nhận của ai bất kỳ điều gì, kể cả là ``đồ thừa''. Nhưng khi con bé nói là từ cô hàng xóm phòng bên, tôi đành im lặng.

Bà ấy là người phụ nữ trung niên sống đơn thân, ít nói, hay giúp đỡ mấy đứa nhỏ trong khu. Ít ra, tôi chưa thấy bà có hành vi gì đáng ngờ. Từng có một lần bà dúi vào tay Suzu túi bánh gạo (hết hạn) mà tôi vẫn còn để nguyên trong ngăn tủ.

``\vie{Cơm thì chị ăn luôn đi ạ. Em để ở đây từ lúc nãy, giờ chỉ cần hâm lại là được,}'' Suzu tháo vát lấy hai cái bát đũa cho hai chúng tôi, khăn vẫn quấn trán.

Tôi cười mệt mỏi: ``\vie{Ừ, cảm ơn em.}''

Đó là bữa tối\dots\ mà đúng hơn là buổi ăn đêm. Cơm nguội, vài lát củ cải muối, một hộp mực khô đóng hộp, và một quả trứng luộc màu nâu nhạt, có lẽ được mua giảm giá. Hộp cá thì\dots\ tôi nghĩ cũng hơi sát hạn rồi. Tôi ngửi thử trước khi ăn. Mùi vẫn không tới nỗi tệ.

``\vie{Em mua mấy thứ này bằng tiền túi à?}'' tôi gắp miếng cá nướng trong hộp, nghiêng đầu hỏi.

``\vie{Hôm kia cô quản lý cho em tiền trả vỏ lon. Em gom được mười mấy cái á!}''

Tôi mím môi, không biết nên thấy thương cho Suzu hay buồn cười nữa.

\ct

Sau bữa ăn, chúng tôi nằm dài trên tấm đệm, bát đĩa thì để nguyên đó để ngày mai chúng tôi dọn. Gió cuối thu ngoài trời lùa qua cửa chớp nghe lành lạnh.

Đèn tắt, cả căn phòng chật hẹp ấy chỉ còn một màu bóng tối, lấp lóe vài ánh đèn đường từ những gian nhà gần đó rọi vào.

Tôi đang tính nhắm mắt ngủ sau một ngày dài thì Suzu lên tiếng: ``\vie{Onee-sama này\dots}''

``\vie{Hửm?}''

``\vie{Khi em lên cấp ba ấy\dots\ em cũng muốn phụ giúp chị bươn chải nữa.}''

Tôi quay đầu lại, nhìn đôi mắt to tròn của con bé trong bóng tối. Đây không phải là lần đầu tiên Suzu cầu xin tôi một điều như vậy. Và như mọi lần, tôi luôn phản đối chuyện này.

Lần này cũng vậy. Tôi lập tức phản bác bằng giọng nói nhẹ nhàng: ``\vie{Không cần. Chừng nào chị còn kiếm ra tiền, em cứ học đi cho tử tế.}''

``\vie{Nhưng mà Onee-sama\dots}''

``\vie{Không có nhưng nhị gì hết,}'' tôi cắt lời. ``\vie{Học xong, có công việc đàng hoàng, rồi kiếm tiền gấp năm, gấp mười chị bây giờ cũng được. Chứ nếu học hành kiểu dở dở ương ương như thế là sẽ bị chôn chân mãi đấy.}''

Suzu khẽ gật đầu, không nói thêm lời nào nữa. Tôi kéo cái chăn mỏng qua người nó, khẽ vỗ vai: ``\vie{Ngủ đi. Mai chị đi làm sớm.}''

Chúng tôi nằm sát nhau trên tấm đệm cũ đã sờn, chăn mỏng không đủ ấm nhưng có lẽ cũng chẳng cần hơn, khi hai chúng tôi có thể ôm lấy nhau. Với hai chị em chúng tôi, chỉ như vậy là đủ.

Tôi đã từng nghĩ lúc đó, rằng: có lẽ mình còn có thể sống kiểu này thêm vài năm nữa. Tôi phải phấn đấu cật lực thêm để chu cấp cho con bé, ít nhất đến khi nó có thể tự lo được.

Tôi đinh ninh với suy nghĩ đó, trước khi thiếp đi, chìm vào giấc ngủ. Lúc đó đã hơn mười một giờ rưỡi đêm.

\begin{center}
    \uline{Tạm ngừng hồi tưởng.}
\end{center}

``Và rồi?'' giọng Kuon chen vào, kéo tôi khỏi dòng suy nghĩ.

Tôi thở dài. Lời nói của cậu ta làm vỡ vụn đoạn yên tĩnh duy nhất mà tôi còn lưu giữ. Tôi biết tôi có thể nổi giận và mắng một trận vì làm hỏng mạch cảm xúc của tôi, nhưng\dots\ tôi không thể.

``\vie{Rồi thì\dots\ khoảng một giờ sau, \emph{thứ đó} đến.}''

``\vie{Thứ đó?}''

``\vie{Lửa. Và không chỉ lửa không đâu. Một cái gì đó còn khủng khiếp hơn cả lửa.}''

\begin{center}
    \uline{Tiếp tục hồi tưởng.}
\end{center}

Tôi không rõ vì sao mình tỉnh dậy. Cảm giác đầu tiên là nóng. Nóng hầm hập, như thể ai đó đang úp một cái chăn dày cộp lên ngực tôi.

Rồi đến mùi. Một thứ mùi mà không ai muốn ngửi thấy lúc nửa đêm: mùi khét nồng nặc, không giống cháy bếp, kèm theo tiếng gì đó nứt răng rắc như gỗ cháy.

Tôi mở mắt, chớp liên tục. Không khí nặng trĩu, và hơi cay bắt đầu len vào mũi, cảm tưởng như ho sặc sụa tới nơi.

Tôi ngồi bật dậy, tim đập mạnh, cảm thấy có điều chẳng lành. Tôi vội lay vai Suzu dậy: ``Suzu\dots\ \textbf{Suzu! Dậy đi!}''

``\vie{Ưm\dots\ gì vậy, Onee-sama\dots?}'' con bé ngái ngủ, dụi mắt. ``\vie{Vẫn còn sớm mà\dots}''

Ngay khi Suzu dứt câu, một tiếng hét thất thanh từ tầng trên vọng xuống: ``\vie{\textbf{Cháy rồi! Cháy rồi! Có ai không! Cứu tôi với!}}''

Và không chỉ có một tiếng. Từ các phòng xung quanh, từ trên xuống dưới, tôi còn nghe được tiếng la hét, tiếng đập cửa, tiếng chân chạy rầm rập trên cầu thang, tiếng trẻ con khóc, tiếng người lớn gào thét gọi nhau. Mùi khói bắt đầu dày đặc, len lỏi qua khe cửa, khiến tôi ho sặc sụa.

Tôi bật dậy, kéo Suzu lên khỏi đệm, vừa hét vừa vơ lấy những gì trước mắt. ``\vie{\textbf{Cháy rồi! Đứng dậy mau!}}'' tôi hét lớn, vơ lấy điện thoại và đèn pin.

Suzu lúc này mới hoàn toàn tỉnh ngủ. Mắt nó tròn xoe, môi run run: ``\vie{Ch-Cháy\dots? Thật á\dots?}''

``\vie{\textbf{Giờ không phải lúc hỏi!! Mau chạy ra khỏi đây đi!!}}''

Con bé còn chưa kịp hiểu chuyện gì, tôi đã lôi con bé ra khỏi phòng, chân trần chạy qua hành lang tối om, bộ quần áo lúc đi làm chưa thay ra. Mở cửa ra, tôi gần như rụng rời: khói dày đặc bao phủ hành lang, và ánh lửa vàng rực phản chiếu trên bức tường đối diện, tiếng những bước chạy rầm ập khắp nơi, khói đen cuồn cuộn tràn vào từng ngóc ngách.

Tôi nắm tay Suzu, vừa ho sặc sụa vừa kéo nó chạy về phía cầu thang.

Gió lùa ngược mang theo khói và sức nóng phả thẳng vào mặt. Một tiếng nổ bụp! từ đâu đó phía dưới khiến cầu thang chấn động. Tôi ngã khuỵu, kéo theo Suzu. Hai đứa lăn một vòng, đập người vào tường.

``\vie{Onee-sama\dots\ em sợ quá\dots\ Cháy thật rồi hả chị!?}'' Suzu hoảng loạn, mắt mở to, mặt tái mét, rơm rớm nước mắt. Con bé vừa mới tỉnh ngủ, giờ khi tình cảnh nguy cấp ở trước mắt đã run rẩy, chân không đứng vững.

Tôi ôm chặt lấy vai nó, cố giữ bình tĩnh dù tim đập thình thịch: ``\vie{Không sao đâu, Suzu! Bám lấy chị, đừng buông ra! Chúng ta sẽ thoát được!}''

Em tôi vừa chạy vừa khóc, miệng lắp bắp: ``\vie{Em không muốn chết đâu\dots\ Onee-sama!! Cứu em với!}''

Tôi ngước nhìn xuống dưới, sức nóng kinh khủng xốc thẳng vào tôi, từng ánh lửa bay vào người tôi, khói đen mù mịt.

Chúng tôi đổi hướng, chạy ngược lên, nhưng vừa tới chiếu nghỉ tầng ba thì thấy lửa đã chặn cầu thang. Ánh lửa bùng lên đỏ rực, khói đen cuồn cuộn, hơi nóng hầm hập như muốn nuốt chửng tất cả.

Suzu đứng sững lại, chân run lẩy bẩy, hét lên trong tuyệt vọng: ``\vie{\textbf{Chị ơi! Lửa chặn rồi! Làm sao bây giờ!? Em không muốn bị kẹt lại đâu!}}''

Tôi siết chặt tay Suzu, mắt đảo liên tục tìm lối thoát, cố gắng không để lộ sự hoảng loạn: ``\vie{Quay lại phòng! Mau lên! Đừng nhìn xuống dưới!}''

Suzu vừa chạy vừa ngoái đầu lại, nước mắt chảy dài: ``\vie{Onee-sama, em sợ quá\dots\ Em không thở được\dots}''

Tôi kéo mạnh tay em, vừa chạy vừa hét: ``\vie{\textbf{Cố lên, Suzu! Bám chặt lấy chị! Không được buông ra! Chúng ta sẽ ổn thôi!}}''

Hai chị em vừa chạy vừa ho, tiếng chân đập loạn trên nền nhà, tiếng lửa nổ lách tách phía sau, tất cả hòa vào nhau thành một cơn ác mộng không lối thoát.

Chúng tôi quay đầu chạy giữa hành lang nóng hầm hập, tim đập như sắp nổ tung, ba chân bốn cẳng chạy về căn phòng của mình. Tôi đẩy cửa phòng, đóng sầm lại, và khóa trái.

``\vie{\textbf{Suzu! Quần áo! Nước!}}'' tôi gào lên.

``\vie{V-Vâng!}''

Cô bé hiểu ý, lao vào góc phòng, nơi tôi thường tập kết quần áo bẩn mà chưa có dịp rảnh rỗi để giặt. Suzu dội nước vào đống áo, ném cho tôi vài cái. Tôi dùng chúng bịt kín khe cửa, chặn gió và khói vào trong, vừa làm vừa run rẩy, tay lóng ngóng vì hoảng loạn.

``\vie{Onee-sama, khói vào nhiều quá!}'' Suzu vừa ho vừa nói, nước mắt lưng tròng.

``\vie{Cố lên, Suzu! Đưa chị thêm áo nữa!}'' tôi tiếp tục thúc gục, giọng khản đặc, mắt cay xè vì khói.

Tôi chạy lại cửa sổ, cố gắng mở nhưng cửa sổ bị chặn song sắt. Tôi gào lên, đập vào song như điên, tuyệt vọng nhìn ra ngoài.

``\vie{Chị ơi, em không thở được\dots}'' Suzu ôm lấy tôi, toàn thân run bần bật.

``\vie{Dùng khăn ướt bịt mũi vào đi! Chúng ta sẽ thoát được khỏi đây thôi!}'' tôi cố trấn an, nhưng chính tôi cũng không tin vào lời mình nói.

Không có lối thoát.

Lửa từ hành lang đã bắt đầu liếm vào khung cửa, làm tường nóng ran, sức nóng khiến tôi giật mình, bật ngược xuống đất. Từng cuộn khói đen cuồn cuộn tràn qua khe, mọi nỗ lực bịt kín đều vô dụng. Tôi cảm thấy cổ họng bỏng rát, mắt cay xè như sắp mù, hơi thở ngắt quãng.

Suzu ho sặc sụa, ôm chặt lấy tôi, giọng lạc đi vì sợ hãi: ``\vie{Onee-sama\dots Em không muốn chết đâu\dots}''

Tôi siết lấy em, nước mắt cũng trào ra: ``\vie{Chị xin lỗi\dots Chị không bảo vệ được em\dots}''

Cơ thể hai chị em tôi co lại trong góc phòng, chờ đợi trong tuyệt vọng khi khói và lửa ngày càng dày đặc. Mọi nỗ lực đều đổ sông đổ bể. Tôi\dots\ tôi không biết làm gì nữa. Cầu cứu? Nhưng điện thoại mất sóng. Đập cửa? Không ai thèm nghe cả. Giữa một bầy người kêu cứu trong vô vọng vì không thể thoát ra như chúng tôi, chắc gì đã tới lượt chúng tôi có thể thoát khỏi đây?

Tôi chỉ biết ôm lấy em mình, cúi đầu và cắn răng chờ đợi, với niềm hy vọng dần mong manh.

Khi cả căn phòng rực sáng bởi ánh lửa tràn vào, tôi đã nghĩ\dots\ mọi thứ đã là dấu chấm hết.

Tôi ôm Suzu chặt tới mức gần như đâm móp cả khung xương của con bé. Chúng tôi rúc vào góc tường, giữa đống áo ẩm và hai cái ba lô cũ. Tôi định bụng, nếu phải chết, ít nhất phải chết cùng nhau.

Nhưng rồi một thứ gì đó khẽ rung. Sàn nhà rung chuyển nhẹ. Ban đầu, tôi tưởng là do ngọn lửa thiêu cháy kết cấu tòa nhà, rằng cả căn phòng này sẽ hoàn toàn đổ sập.

*ẦM ẦM*

Nhưng hóa ra, lại không phải vậy.

Giữa bức tường đối diện cửa sổ, không gian rách toạc.

Một lỗ xoáy xuất hiện, xoay tròn như hố đen, với tia sáng xanh tím chạy dọc biên viền, tạo ra âm thanh kỳ dị như tiếng gió rít pha trộn với tiếng kim loại ma sát. Cảm giác như nhìn vào thứ không nên tồn tại trong thế giới này.

``\vie{Cái gì\dots\ cái gì thế kia, Onee-sama\dots!?}'' Suzu thốt lên, vừa ho vừa kéo tôi lại gần hơn.

``\vie{C-Chị cũng không biết!}'' tôi thảng thốt, ôm chầm lấy cô em gái mình.

Một lực hút bật lên, kéo mọi thứ xung quanh: giấy tờ, quần áo, cả chiếc đệm chúng tôi đang ngồi lên. Tất cả bay về phía cái lỗ xoáy không mời mà tới.

Tôi cảm thấy người mình nhẹ bẫng, như bị giật mạnh khỏi mặt đất.

Và rồi\dots\ mọi thứ tối sầm lại chỉ trong nháy mắt.

\begin{center}
    \uline{Kết thúc hồi tưởng.}
\end{center}

``Lỗ xoáy đó\dots\ trông như thế nào?'' Kuon bất ngờ giương khuôn mặt nghiêm nghị.

Tôi chớp mắt, chậm rãi gật đầu, cố nhớ chi tiết về cái lỗ xoáy kia. ``\vie{Xem nào\dots\ Nó tròn, xoáy ngược chiều kim đồng hồ, viền màu xanh tím\dots\ trôi nổi giữa không trung. Giống như một vết rách thực tại.}''

Cậu ta khoanh tay, nắm mắt lại đầy suy tư, rồi đáp lại không một chút chần chừ: ``\vie{Tôi biết nó. Tôi cũng từng thấy nó. Giữa ngã tư gần khu trung tâm. Thứ đó đã đưa tôi tới đây.}''

Tôi giật mình, mắt mở to vì ngạc nhiên. ``\vie{Cậu cũng thấy nó ư? Vậy\dots\ cậu có biết nó tới từ đâu không?}''

Kuon thoáng lúng túng, gãi đầu, mắt hơi đảo ra chỗ khác như tránh né ánh nhìn của tôi. ``\vie{\dots Nó cũng hơi khó giải thích, và cũng loằng ngoằng nữa. Nhưng tóm lại, tôi biết.}''

Tôi định hỏi thêm, nhưng Kuon lập tức chuyển chủ đề, như muốn né tránh: ``\vie{Sau đó thì sao?}'' cậu hỏi tiếp.

Tôi rướn mắt nhìn lên trần của hư không. Cổ họng khô đắng. ``\vie{Khi tỉnh lại, tôi thấy mình nằm trên  cái đệm cháy xém. Lúc đó tôi còn chưa hiểu chuyện gì vừa xảy ra nữa.}''

``\vie{Chính là cái mà em nhắc đó!}'' Suzu chen vào, mắt sáng lên.

Tôi gật nhẹ. ``\vie{Nhưng rồi, tim tôi thắt lại. Tôi ngồi bật dậy, đưa mắt nhìn quanh. Căn phòng cháy biến mất. Khói và lửa không còn. Mọi thứ trở nên đen đặc, không gian vô định chính là nơi mà chúng ta đang đứng bây giờ.}''

``\vie{Nhưng Suzu\dots\ không ở đó,}'' tôi nói, giọng vẫn còn run run, khi tôi buộc phải nhớ lại trải nghiệm kinh hoàng ấy. ``\vie{Tim tôi như rơi vào hố băng.}''

``\vie{Tôi gào gọi, chạy loạn trong hư không,}'' giọng tôi thấp trầm xuống. ``\vie{Mỗi bước chân đều như vô nghĩa trong khoảng trống mịt mùng. Tôi hét đến khản cổ\dots\ nhưng không một tiếng đáp lại. Và cứ thế, tôi vừa chạy, vừa gọi trong vô vọng, lang thang mãi ở nơi quái đản này.}''

Vào lúc đó, tôi tưởng tôi đã mất con bé rồi. Tôi còn nghĩ tới trường hợp xấu hơn, rằng Suzu đang ở một nơi hiu quạnh nào đó mà tôi không biết. Nếu tôi không tìm thấy con bé, không biết tôi sẽ phát rồ tới mức nào.

Thậm chí\dots\ nếu con bé có mệnh hệ gì, tôi sẽ dằn vặt bản thân mình tới suốt đời. Có khi tôi sẽ sống như một cái xác biết đi trong cái thế giới trống rỗng này, không mục đích, không hy vọng. Bơ vơ\dots\ cô độc\dots\ mãi mãi.

Nhưng ơn trời, nó đã không xảy ra. Suzu đã trở về an toàn, không một vết trầy xước-- ờm, không hẳn, nhưng đúng hơn là con bé không bị thương tích gì nghiêm trọng. Con bé vẫn cười được, vẫn ôm lấy tôi, vẫn gọi hai tiếng ``Onee-sama''\dots

Nhưng điều khiến tôi sửng sốt hơn\dots

\dots là nó không trở về một mình. Suzu được một người con trai - một người lạ mặt, lạc vào nơi này - giúp đỡ, chăm sóc, dắt đi, và cuối cùng đưa nó trở lại bên tôi. Cậu ta không hề đòi hỏi gì. Không hỏi xin sự công nhận. Không đợi lời cảm ơn. Thứ duy nhất mà cậu ta mong muốn chỉ là thoát khỏi nơi chết tiệt này càng sớm càng tốt.

Tôi đã nhìn thấy nhiều người đàn ông, từ khi còn nhỏ, cho đến tận lúc lớn lên và đi làm thêm. Phần lớn trong số đó\dots\ đều khiến tôi cảm thấy khó chịu, mệt mỏi, hoặc tệ hơn, gớm ghiếc. Ánh mắt bủa vây. Lời nói hai nghĩa. Tay chân vô ý thức.

Nhưng cậu ta thì khác.

Cậu ấy - Hikawa Kuon - có thể hơi thô lỗ, hơi bất cần, và rõ ràng là không thân thiện gì mấy, nhưng\dots\ không có ham muốn nào trong mắt cậu ta cả. Cậu ta nhìn tôi, nhìn Suzu\dots\ như nhìn những con người thực sự, không tô vẽ, không nuốt chửng, không tỏ vẻ dâm dục gì.

Tôi không còn nhớ lần cuối mình được thấy nó là khi nào nữa, kể từ khi bố tôi qua đời.

``\vie{\dots Em đáp xuống ở đâu đó khác,}'' Suzu nói nhỏ. Tôi quay sang con bé. Khuôn mặt của cô bé hiện một chút đo đỏ.

``\vie{Có cái xấp quần áo bẩn mà em giặt dang dở chiều hôm đó. Em đáp đúng vào đống đó nên cũng không đau lắm,}'' con bé thó thé, khẽ run run.

``\vie{Nhưng xung quanh em\dots\ chỉ có bóng tối. Không có tiếng động, không có ai.}'' Đôi mắt con bé trực trào nước mắt, như thể câu chuyện Suzu kể mới chỉ ngày hôm qua vậy. ``\vie{Em đã gọi chị, gọi Onee-sama thật lớn, nhưng vẫn có hồi đáp. Em khóc, em chạy\dots\ mà chỉ thấy toàn là những thứ đen sì, kỳ dị. Em tưởng chị đã bị\dots}'' con bé nghẹn lại, không nói được hết câu.

Tôi siết chặt tay lại, tự cảm thấy lòng mình đau nhói. Suzu hẳn đã phải chịu đựng cảm giác lạnh lẽo ấy rất lâu rồi. Giá như tôi có thể ở đó, ôm chặt con bé vào lòng, rồi xoa nhẹ lưng vừa nói: ``Không sao đâu, Suzu\dots\ Có Onee-sama ở đây rồi.'' Tôi ước rằng tôi có thể sát cánh bên em mọi lúc.

Dù vậy, điều quan trọng, là con bé vẫn khỏe mạnh, vẫn có thể đứng đây nói chuyện với tôi. Vậy là quá đủ.

\dots

\dots

Tiếng thở dài của Kuon làm cho bọn tôi giật mình. Cậu khoanh tay, gục đầu xuống, bản mặt nửa chán nản, nửa đồng cảm, rồi bất giác buông ra một câu nói khiến cả hai chị em tôi chết lặng vài giây: ``\vie{Nếu không vì cái lỗ xoáy kia, chắc hai người cũng đã lên thiên đàng tề tựu với bố mẹ rồi nhỉ?}''

Một sự im lặng bao trùm. Điều đó\dots\ đúng đến đau lòng. Một lời nói trực diện không hề xuôi tai, nhưng sự thật là như vậy.

\pov{Hikawa Kuon}

Kaigou khẽ bật cười với cái kiểu nói thẳng đến trần trụi của tôi. Không hẳn là tiếng cười vui, cũng chẳng hẳn là giễu cợt. Nó giống như tiếng thở ra khi một người quá mệt để tức giận thêm.

``\vie{Nếu có thể quay về nhà thật\dots}'' cô ấy nhìn về phía bóng tối xa xăm, ``\vie{chắc tôi vẫn sẽ tiếp tục làm thêm để nuôi Suzu ăn học. Vẫn sẽ là cái siêu thị cũ đó, hay bất kỳ việc gì đủ tiền trả tiền trọ. Chắc sẽ vẫn là căn phòng đó, hoặc nếu cháy rồi thì xin dọn vào một chỗ nào còn tệ hơn\dots\ Tôi quen rồi.}''

Cô ngừng một nhịp. Nhìn sang Suzuki-chan đang lắng nghe bên cạnh.

``\vie{Nhưng tôi cũng muốn\dots\ thay đổi. Tôi muốn Suzu có một cuộc sống khác. Một thứ gì đó đàng hoàng hơn. Một nơi mà con bé không cần phải cúi đầu với bất kỳ ai chỉ vì nó là con gái. Không cần phải sống trong cảnh nơm nớp, run rẩy mỗi khi nghe tiếng bước chân ngoài hành lang. Còn tôi, chắc thế nào cũng được.}''

Tôi im lặng lắng nghe lời thỉnh cầu nhỏ nhoi nhưng thành thật ấy. Trong đầu tôi, rất nhiều hình ảnh chạy ngang qua cùng một lúc:

Suzuki-chan với ánh mắt khẩn cầu yếu ớt khi lần đầu cầu xin tôi đừng trách chị mình.

Kaigou lạnh lùng và dữ dằn, nhưng cũng là người đã thẫn thờ khi không thấy em gái đâu.

Hai người họ, cùng tồn tại với nhau trong một thế giới nơi đàn ông chiếm đa số, nơi quyền lực, tiền bạc, và sự bạo lực không chia đều cho phái yếu.

Những buổi tối dùng cơm thừa, đồ hết hạn. Căn phòng trọ bé xíu với tiếng cãi vã từ tầng dưới.

Và hơn hết, là hai nụ cười hiếm hoi khi họ kể lại mẩu chuyện hài hước về quạt trần, hay khi nhắc đến bộ truyện tranh họ cùng đọc dưới ánh đèn vàng vọt yếu ớt.

Tôi thở nhẹ. Trong tôi có một cảm giác gì đó\dots\ không hẳn là thương hại, cũng chẳng phải cảm thông kiểu sáo rỗng.

Chỉ là, tôi hiểu cảm giác của họ.

Một người có xuất thân là gia đình trung lưu có thu nhập ổn định, vẫn còn có cả bố mẹ, còn có họ hàng giúp sức, vẫn có của ăn của mặc đàng hoàng như tôi, nhưng tôi hiểu được những gì họ đã từng trải qua, dù không hoàn toàn.

Tôi đã từng trải nghiệm sống ở một nơi tồi tàn rách nát, đã từng ngủ trong môi trường khắc nghiệt của tiết trời Kawachi khi hè tới hay khi đông sang ở bên xưởng bố tôi mỗi khi xuống phụ việc. Biết cảm giác mà muỗi bu đầy người khi đang ăn mì tôm giữa cái hè oi ả, dù rằng đây chỉ là trải nghiệm ngắn hạn và có thể rất khác so với những gì mà hai chị em đã từng trải qua.

Và tôi biết mình cần phải làm gì. Một quyết định có phần hơi bốc đồng, nhưng cần thiết.

``Này.''

Kaigou quay lại, đôi mắt cô ấy vẫn cảnh giác, nhưng ít gai hơn ban đầu. ``\vie{Nếu hai người quay về thế giới cũ, cô sẽ lại tiếp tục sống như trước, phải không?}''

``\vie{Có lẽ vậy. Còn gì khác để làm đâu?}'' cô nàng nhún vai.

``\vie{Nhưng nếu tôi nói\dots\ tôi có cách khác thì sao?}''

Suzuki ngẩng đầu lên. Kaigou cau mày. Dĩ nhiên là không một ai trong số hai người họ tin vào điều tôi nói rồi.

Tôi liếc sang hai người. ``\vie{Tôi sẽ đưa hai người đến chỗ tôi. Không phải thế giới này. Một nơi khác hẳn, nơi mà thời tiết có thể khắc nghiệt không kém gì nơi mà hai người từng sống, nhưng lòng người thì mát hơn nhiều.}''

Tôi nhếch môi. ``\vie{Ở đó, phụ nữ được tôn trọng. Hàng xóm sẵn sàng nhường cơm sẻ áo, tối lửa tắt đèn còn có nhau. Và hơn hết, hai người sẽ có một mái nhà cùng với gia đình tôi. Một nơi để bắt đầu lại từ con số không.}''

Tôi dừng lại, quan sát sắc mặt Kaigou. Suzuki thì đã mở to mắt, gần như không tin vào tai mình. Những gì tôi nói có lẽ không khác nào một vùng đất thần tiên mà có lẽ chỉ có thể nằm mơ mới thấy.

Ngoại trừ việc đó không phải mơ. Đó là sự thật.

``\vie{Tôi không dám chắc nơi tôi sống sẽ là thiên đường, nhưng tôi hứa hẹn rằng ít nhất đó không phải là một nơi địa ngục, không phải là tù đày. Đó sẽ là nơi mà hai người có thể sống thật với bản thân nhất, có thể sống theo cách mình muốn.}''

Kaigou nheo mắt. Sự im lặng giữa ba chúng tôi dường như đặc quánh lại. Tôi biết cô ấy sắp nói gì đó. Có thể là từ chối, có thể là chất vấn.

Tôi không chắc. Nhưng tôi sẵn sàng để nhận lời hồi đáp từ cô ấy.

``\vie{Cậu nói\dots\ muốn đưa chúng tôi đến `nơi tốt hơn' sao?}'' giọng cô vang lên. Nhẹ nhưng sắc như kim. Tôi vẫn cảm thấy có một sự đề phòng rất cao, nhưng tôi không ngạc nhiên.

Tôi gật đầu, không hề tránh né. ``\vie{Đúng, là quê của tôi. Nơi tôi lớn lên. Một căn nhà nhỏ chưa tới ba chục mét vuông, nhưng đủ ấm, đủ để ngồi học, để ăn, để ngủ mà không sợ trộm hay chuột.}''

Tôi liếc vào chiếc đồng hồ đeo tay mà bố tôi đã trao cho tôi, dù hơi sờn cũ nhưng vẫn sử dụng được tốt.

``\vie{Ở đó, không ai quan tâm tôi mặc gì, ăn gì, hay đi làm lúc mấy giờ. Chỉ cần không làm phiền bà nội và mấy bác dưới nhà thì muốn làm gì cũng được.}''

Kaigou khoanh tay. Ánh mắt vẫn cảnh giác. ``\vie{Nghe giống như một nơi thần tiên nhỉ,}'' cô nàng cất giọng mỉa mai nhẹ. ``\vie{Nhưng lũ đàn ông xung quanh tôi cũng từng ngọt ngào như thế, trước khi-}''

``\vie{Tôi thì không bao giờ bịa chuyện công việc đâu,}'' tôi ngắt lời, không cần lớn tiếng. ``\vie{Tôi hiện tại là tổng quản lý tại một công ty giải trí lớn ở Nhật, kiêm luôn mấy công việc lặt vặt ở đó. Cũng chính công việc ấy giúp tôi đủ tiền gửi học phí cho mấy đứa sinh viên tôi nhận bảo trợ, chăm sóc cả em gái tôi, và cả tiền để nuôi thêm hai `con vịt giời' nữa, nếu cần. Chưa kể bố mẹ tôi vẫn còn đang đi làm, nên cũng phụ giúp được phần nào.}''

Cô ấy nheo mắt, định lớn tiếng. ``\vie{Cậu nghĩ chỉ vì có tiền, có công việc, thì có thể mua lòng thương hại sao?}''

Tôi khẽ cười. Không phải kiểu sắc sảo, mà là chua chát với một cô gái đang mất niềm tin vào xã hội. ``\vie{Tôi không hề thương hại. Tôi đã từng nằm co trong phòng trọ ẩm mốc đầy muỗi, đi học trong bộ đồng phục cũ sờn, trải qua một khoảng thời gian mà cả gia đình tôi chỉ có bánh mì để cầm cự cả tháng, trong khi nhìn người ta có bữa cơm thịt hẳn hoi đấy. Tôi biết cái cảm giác đó. Cái cảm giác mà chị em cô đã và đang trải qua.}''

Tôi nhìn thẳng vào mắt Kaigou. ``\vie{Đó không phải là thương hại. Đúng hơn, là thấu hiểu. Mà thấu hiểu thì khác. Rất khác.}''

Dù nói đến tận đây, cô nàng vẫn chưa hề buông lỏng. Cô tiếp tục phản bác: ``\vie{Nhưng thế giới của cậu dù có đẹp đến đâu\dots\ vẫn có nguy cơ. Tôi biết rõ mà.}''

Tôi gật đầu. Không hề phủ nhận. ``\vie{Đúng là như thế. Có thể, vẫn có người sẽ dòm ngó. Có thể, vẫn sẽ có những tên dâm đãng ở đâu đó ngoài kia. Nhưng\dots}''

Tôi ngẩng đầu nhìn lên khoảng trống đen thẫm trên kia, nơi lẽ ra là bầu trời. ``\vie{Nếu cô hét lên ở Kawachi tôi sống, người ta sẽ nghe thấy. Và họ sẽ chạy đến. Có khi là bác tổ trưởng dân phố, có khi là bà bán bún, có khi là mấy cậu sinh viên tầng trên. Họ không cần biết cô là ai, nhưng họ sẽ ra tay giúp, và không đòi hỏi gì cả. Thậm chí còn cho đi là đằng khác.}''

Kaigou mở miệng định nói gì đó. Tôi không để cô kịp chen vào. ``\vie{Bởi vì ở đó, người ta vẫn còn giữ được cái gọi là `tình người.' Tôi không biết hoàn toàn thế giới của cô vận hành như thế nào, nhưng tôi có thể đảm bảo rằng thế giới nơi tôi ở tử tế hơn nhiều.}''

Một sự im lặng bao trùm. Tôi thấy cô nàng đang đấu tranh tư tưởng.

Tôi lập tức hạ giọng, chậm rãi như một người bố đang giảng dạy con cái: ``\vie{Không phải đàn ông nào cũng xấu. Không phải ai cũng là ác quỷ trong hình hài con người. Cô nên mở rộng tấm lòng mình một chút đi, thế giới không phải lúc nào cũng là một gam màu tối xám xịt đâu.}''

Tôi quay sang Suzuki-chan, người từ nãy vẫn lặng im nghe.

Cô bé run lên. Nhưng ánh mắt của em ấy trông không có vẻ gì là sợ sệt như lúc mới gặp cả. Mà một cái gì đó như là sự quyết tâm.

``\vie{Em\dots\ em cũng muốn một nơi như thế, Onee-sama,}'' giọng của Suzuki-chan như nén chặt, mắt đỏ hoe vì xúc động. ``\vie{Một nơi mà bọn em không phải sống rón rén, không phải nghe tiếng khóa cửa là hoảng loạn. Một nơi mà em\dots\ có thể nấu ăn, học bài và\dots\ được ngủ ngon mà không phải nghĩ ngợi.}''

Tôi thấy khóe môi của người chị gái giật giật. Dám chắc cô đang nghĩ: ``\textit{Không phải cả Suzu nữa chứ!? Tại sao em lại về phe cậu ta!?}'' nhưng tôi quyết định không nói ra.

``\vie{Em\dots\ Em tin Kuon-nii-san,}'' cô bé nói tiếp, lần này dõng dạc.  ``\vie{Anh ấy\dots\ chưa từng đòi gì ở em. Chỉ giúp em, dù chẳng quen biết gì cả. Nếu là anh ấy\dots\ em muốn sống ở thế giới của anh.}''

Kaigou không đáp ngay. Tôi thấy ánh mắt cô ấy lạc đi một chút. Giống như khi ai đó nhìn thấy một con đường cũ, nơi họ từng đi qua.

Tôi biết ánh nhìn đó. Vì tôi cũng từng đứng ở ngã ba ấy trước khi tôi điền đơn tuyển dụng cho \hll\ - giữa việc ăn bám bố mẹ và trở thành một kẻ bất tài và việc phải thay đổi bản thân để tự nuôi sống bản thân. Tôi biết, cách ví von đó không đúng lắm với tình cảnh hiện tại mà người chị đang gặp phải, nhưng về cơ bản là cũng gần như thế.

\pov{Hikawa Kaigou}

``\eng{Now is your choice,}'' giọng cậu ta không cao, cũng chẳng lạnh lùng, chỉ lặng như mặt hồ đêm tối. ``\eng{You can go back to your world and continue living a miserable life full of vengeance and hatred for, I don't know, five to ten years, and even for the rest of your life, or come with me and start over from stratch for a brand new life with my family, living in a world with full of hope, respect, and chances.}\ \vie{Tôi dám chắc là cô cũng biết rõ một chút tiếng Anh cơ bản để hiểu những gì tôi nói chứ?}''

Tôi gật đầu, mắt vẫn chăm chăm nhìn vào cậu.

Tôi hiểu. Một cách hoàn hảo.

Những từ đó như nhát dao cắt vào chính ký ức tôi.

Cậu ta vẫn đứng đó. Không thúc ép, không cầu xin, không đưa tay ra kéo tôi theo, mà lại rất đủng đỉnh, kiểu ``chọn gì cũng được, tôi vẫn sẽ tôn trọng quyết định của cô.'' Cậu chỉ đứng đó như một điểm sáng duy nhất trong cái không gian âm u này, nhìn tôi bằng ánh mắt mà tôi không quen thuộc, một ánh mắt của sự tin tưởng không điều kiện.

Suzu thì sao?

Con bé đã chọn rồi. Không cần nói nhiều, không cần cân nhắc. Với nó, có lẽ đơn giản là: ``Anh ấy giúp em, nên em tin.'' Vậy thôi. Tấm chiếu chưa trải thì thường dễ tin người là vậy. Nhưng, liệu tôi có thể chỉ vì con bé tin, mà dễ dàng chấp nhận một cuộc đời hoàn toàn mới?

Nếu là tôi của một năm trước, không, chỉ cần một tuần trước, thậm chí là một ngày trước, chắc chắn tôi sẽ khinh thường những lời như thế. ``Gia đình'', ``khởi đầu mới'', ``tình người'', toàn là những cụm từ tô vẽ cho những cạm bẫy đẫm mùi đường mật.

Nhưng bây giờ, khi đứng đây, giữa lằn ranh mong manh của tất cả những gì tôi đã mất và những gì tôi có thể có, tôi mới nhận ra một điều.

Tôi đã quá mệt mỏi rồi.

Tôi mệt với việc sống như thể cả thế giới nợ tôi điều gì đó. Mệt với việc lúc nào cũng phòng thủ, mệt với cái cách mỗi ngày tôi phải gồng mình chống lại đàn ông, xã hội, người thân, cả chính nỗi sợ trong tôi.

Tôi mệt với một xã hội nơi mà phụ nữ bị đe nẹt một cách không thương tiếc, lúc nào cũng đè đầu cưỡi cổ tôi, lúc nào cũng đặt phái yếu vào thế nguy hiểm.

Tôi từng nói với Suzu: ``Chị sẽ bảo vệ em khỏi lũ người gớm ghiếc kia'' không biết bao nhiêu lần. Nhưng rồi\dots\ có phải tôi cũng đang níu lấy con bé để không bị ngã quỵ? Có phải tôi sống đến giờ chỉ vì con bé vẫn còn bên tôi?

Tôi chưa từng dám mơ đến việc được sống bình thường kể từ khi hai chúng tôi chỉ còn lại nhau. Chỉ cần một cuộc đời không ai động chạm tôi, không ai khinh rẻ tôi, không ai nhìn tôi như một món hàng đã là quá xa xỉ.

Mà cái cậu kia, cái người chẳng biết tôi là ai, chẳng nợ tôi điều gì, lại dám nói rằng hãy đến với cậu ta và bắt đầu một cuộc sống mới với gia đình của cậu ư?

Gia đình?

Tôi thậm chí còn quên mất cảm giác ấy như thế nào. Đã bảy năm rồi kể từ khi tôi được cảm nhận lần cuối.

Tôi quay sang Kuon. Gã trai trẻ với dáng vẻ chán đời, tóc tai như tổ quạ, ăn nói lỗ mãng ngang ngược, thái độ phớt lờ chuẩn mực kia lại chính là người khiến con bé mỉm cười trở lại, tỏ vẻ hoàn toàn thoải mái mà không hề lo sợ điều gì cả.

Tôi nhìn Suzu. Ánh mắt con bé vẫn hướng về tôi, nhẹ nhàng mà tha thiết. Không phải là ánh nhìn chờ đợi hay ngây thơ như mọi khi, mà đó là ánh nhìn đã biết chọn lựa. Một ánh nhìn khiến tim tôi thắt lại.

Con bé tiến lại gần, khẽ nắm lấy vạt áo tôi. Ngập ngừng, nhưng dứt khoát: ``\vie{Em không muốn quay lại nơi đó nữa\dots\ Nếu chị không đi, em vẫn sẽ đi theo anh ấy. Nhưng\dots\ em muốn được đi cùng chị cơ, Onee-sama.}''

Câu nói đó khiến tôi chết lặng.

Suzu chưa từng nói câu nào như vậy. Từ bé đến giờ, lúc nào nó cũng chỉ răm rắp nghe lời tôi. Chưa bao giờ đi ngược lại tôi. Chưa bao giờ chọn ai khác. Tôi là thế giới duy nhất của nó, và nó là lý do duy nhất khiến tôi tiếp tục bước đi.

Vậy mà giờ đây, con bé lại có thể rời khỏi tôi mà đi theo một người đàn ông khác, một người mà tôi chỉ mới vừa chấp nhận chưa đầy một ngày.

Tôi cắn môi. Một cơn co thắt chạy dọc lồng ngực.

``Em vẫn sẽ đi\dots\ nhưng em muốn được đi cùng chị.''

\dots Không lẽ tôi đang mất dần con bé về phía gã ta rồi sao? Tôi\dots\ không lẽ sẽ bị bỏ rơi?

Tôi thấy lòng mình chao đảo. Tựa như mặt đất vừa sụp dưới chân. Tôi chưa từng cho phép mình yếu đuối. Nhưng giờ đây, tôi không còn che giấu được nữa. Không còn mạnh mẽ được nữa.

Tôi buột miệng: ``\vie{Chị mệt quá, Suzu à\dots\ Mệt lắm rồi\dots}''

Tôi không biết tôi vừa nói bằng giọng gì. Có thể là khản đặc, có thể là ứ đọng. Chỉ biết là cánh tay tôi run lên. Đầu tôi cúi gập. Lòng ngực nấc lên một tiếng vô thức. Tôi không biết mình đã rơi nước mắt từ khi nào.

Ngay khi tôi ngỡ mình sẽ ngã quỵ, một bóng người bước tới.

Một vòng tay nhẹ nhàng, vừa đủ ấm, vừa đủ chặt, ôm lấy tôi từ phía trước.

Không có dục vọng. Không có toan tính hiểm độc. Không lời an ủi sướt mướt.

Chỉ là một cái ôm, giống như ngày xưa mẹ tôi từng ôm tôi khi tôi bị ngã, khi tôi gặp ác mộng, khi tôi khóc vì nhớ bố. Một cái ôm khiến tôi không thể kìm được nữa.

``\vie{Cô đã phải chịu đựng quá nhiều rồi\dots}'' Kuon nói, giọng trầm, như đang nói với chính tim mình. ``\vie{Cứ khóc đi. Đừng kìm nén nữa. Người mạnh mẽ đôi lúc cũng có quyền được yếu đuối mà.}''

Chỉ chờ có thế, tôi bật khóc.

Một tiếng khóc òa giữa một không gian tối tăm và trống rỗng này, thứ tiếng khóc nức nở, run rẩy, đứt đoạn, nghẹn ứ ở cổ, thứ tiếng khóc mà tôi đã chôn chặt suốt bảy năm trời. Cơ thể tôi co lại như một đứa trẻ. Tay bấu vào lưng áo Kuon. Mắt nhòe nước. Tôi không thấy gì nữa, không nghe gì nữa ngoài tiếng thổn thức phát ra từ chính ngực mình.

Suzu cũng vòng tay ôm tôi từ bên cạnh, nhẹ nhàng dụi đầu vào vai tôi. Trong khoảnh khắc đó, tôi không còn là Kaigou, một người gồng gánh, một kẻ mạnh mẽ bảo vệ cô em gái thân yêu, mà chỉ đơn giản là một đứa con gái bình thường, đang được hai người duy nhất còn lại trên thế giới ôm lấy.

Tôi rấm rứt khóc như một đứa bé, khóc đến cạn hơi, khóc đến khản cổ.

Giữa tiếng nấc, tôi đấm nhẹ vào vai cậu ta một cái, vừa đủ để không khiến cậu đau mũi vẫn sụt sịt: ``\vie{Đồ ngốc\dots\ sao lại tốt bụng như vậy chứ\dots?}''

Kuon cười khẽ, vẫn giữ tay trên lưng tôi: ``\vie{Vì tôi thích thấy cô khóc hơn là thấy cô chết chìm trong đau khổ.}''

Tôi chẳng biết nên đáp gì, mà chỉ tiếp tục trút hết nỗi lòng của mình qua những cú đấm vào lưng, vào vai, qua những hàng nước mắt.

Chỉ biết là lần đầu tiên sau nhiều năm, tôi đã cảm thấy nhẹ nhõm. Như thể vừa trút được cả một tảng đá khỏi tim mình.

\pov{Hikawa Suzuki}

Tôi đứng nép bên cạnh, vòng tay ôm lấy chị, đầu tựa nhẹ vào bờ vai gầy guộc mà tôi đã quen dựa vào suốt mười năm nay. Ngực tôi phập phồng, cổ họng nghèn nghẹn, nhưng là vì một cảm giác lạ lắm: một sự ấm áp, nhẹ nhõm và biết ơn.

Tôi chưa từng thấy Onee-sama như thế trong đời.

Từ nhỏ đến lớn, Onee-sama luôn là người đứng trước, luôn là bức tường vững chãi che chắn cho tôi khỏi giông tố. Dù bị đời vùi dập đến đâu, chị cũng không than thở một lời. Dù có kiệt sức mấy, chị vẫn sẽ gồng lên cười. Dù phải ăn mì gói hết tháng, chị cũng sẽ nhường tôi phần ngon hơn. Tôi chưa từng thấy chị khóc. Thậm chí, có những lúc tôi tự hỏi: ``Chị mình có thực sự biết đau khổ không?''

Vậy mà giờ đây, chị đang òa khóc trong vòng tay của một người con trai\dots\ không, không phải là bất cứ người con trai nào, mà là Kuon-nii-san.

Chính anh ấy đã khiến chị tôi sụp đổ, để rồi được nâng lên.

Chính anh ấy đã nói những lời mà tôi tin rằng không ai khác dám nói.

Chính anh ấy, người đã cứu tôi, đã kéo tôi khỏi vực thẳm cô độc trong thế giới này, và giờ\dots\ đang mở ra một lối thoát cho cả hai chị em tôi.

Tôi chưa từng tin tưởng ai ngoài chị mình. Nhưng giây phút này, khi thấy nụ cười đầy nước mắt của chị, một nụ cười thật sự, không gượng gạo, không giả vờ, không che giấu, tôi biết mình đã đúng khi chọn tin vào Kuon-nii-san.

Đó là lần đầu tiên sau nhiều năm, tôi thấy Onee-sama được là chính mình. Không phải người chị cả, không phải người chống đỡ cả gia đình, mà chỉ đơn thuần là một cô gái, đang được yêu thương và được phép yếu đuối.

Trái tim tôi chợt co thắt. Nhưng lần này không phải vì đau khổ, mà vì xúc động.

Nếu được đi theo Kuon-nii-san, nếu được sống trong thế giới mà anh ấy miêu tả, một nơi mà phụ nữ không bị coi thường, nơi người ta sẽ đứng lên bảo vệ nhau thay vì ngoảnh mặt làm ngơ, thì tôi tin, chị tôi sẽ được sống đúng nghĩa.

Và tôi sẽ được thấy chị cười nhiều hơn.

Lần đầu tiên trong đời, tôi thấy một tương lai không còn u ám.

Một tương lai có Onee-sama và có anh ấy. Một tương lai xán lạn hơn đang chờ đón hai chị em chúng tôi.

\ct

\pov{Hikawa Kuon}

``\vie{\dots Suzu đi đâu thì tôi sẽ đi đó.}''

Đó là một cách khác để bày tỏ rằng ``tôi sẽ đi cùng với cậu, Kuon,'' nhưng cô nàng không chịu nói ra.

Tôi ngẩng lên nhìn Kaigou, nheo mắt lại. Cô ấy nói câu đó bằng cái giọng nửa như lạnh nhạt, nửa như giấu điều gì trong cổ họng. Nhưng ánh mắt thì lại không hề giấu nổi.

Đôi mắt của niềm hy vọng như một người cuối cùng đã tìm ra được bến đỗ của mình sau một hành trình mệt nhoài.

Tôi cười nhẹ, chậm rãi hỏi lại: ``\vie{Cô chắc với lựa chọn đó chứ?}''

Kaigou lập tức phồng má lên phụng phịu. Không còn vẻ hầm hố hay đề phòng như trước nữa, cô thậm chí còn quay mặt đi, gắt khẽ: ``\vie{\dots Cậu mà nói thêm câu nào kiểu thế nữa là tôi đổi ý bây giờ đấy.}''

Tôi nhướng mày: ``\vie{Nhưng cô biết là cô không thể, đúng chứ?}''

Lập tức, Suzuki-chan phá ra cười, tiếng cười lanh lảnh như chuông ngân, còn Kaigou thì\dots\ đỏ mặt. Rất đỏ. Má vẫn phồng, miệng thì lầm bầm: ``\vie{Thế mà vừa nãy cậu còn nói lời yêu thương lắm\dots\ Quá đáng\dots}''

Tôi nhún vai, bật cười. Không rõ tại sao, nhưng phản ứng có phần hơi trẻ con, hơi ương ngạnh ấy làm tôi thấy lòng nhẹ đi rất nhiều. Nhẹ hơn cả khi hoàn thành một hợp đồng lớn hay khi đẩy được một tài năng lên đỉnh những người thịnh hành nhất.

Lúc này, trong ánh sáng lập lòe tím nhạt mơ hồ của The Void, Kaigou không còn là người chị cả gồng mình gánh vác cả thế giới nữa. Suzuki cũng không còn là cô bé mồ côi co ro trong góc tối.

Cả hai đã chọn một con đường mới.

Và tôi sẽ là người dẫn lối.

Tôi biết việc này sẽ không hề dễ. Họ sẽ cần thời gian để làm quen, để hòa nhập, để hiểu rằng thế giới mà tôi đưa họ tới không phải thiên đường, nhưng là một nơi tốt hơn, nhân văn hơn.

Tôi sẽ làm mọi thứ trong khả năng của mình.

Vì một mái ấm mà họ xứng đáng có được.

Vì một tương lai mà họ chưa từng dám mơ.

Vì với tôi, Kaigou và Suzuki-chan ở thế giới cũ đã chết rồi.

Giờ, họ sẽ là những con người mới. Trong một cuộc đời mới, ở bên cạnh tôi, với gia đình tôi, và những người mà tôi tin tưởng nhất.

Tôi đã hứa là tôi sẽ làm, vì quân tử không bao giờ nói hai lời. Truyền thống của gia đình tôi là: chỉ có thể hứa những điều mà tôi có thể làm được, và tôi sẽ cố gắng hết sức để thực hiện điều đó.

\dots Nhưng đấy là nếu chúng tôi có thể rời khỏi nơi khỉ ho cò gáy này. Đây mới chính là nhiệm vụ cốt yếu nhất.

\ct

Trong lúc tôi còn đang nghĩ xem tôi sẽ giải quyết vấn đề trước mắt với hai cô gái này như thế nào, thì bỗng nhiên một tiếng gọi lớn kéo tôi về thực tại: ``\textbf{Onii-sama!}''

Tôi khựng lại, mắt mở lớn. Kaigou cũng tròn mắt ngạc nhiên. Phải mất một nhịp, tôi mới hiểu được người vừa lên tiếng chính là Suzuki-chan. Cô bé vừa mới gọi tôi bằng cách gọi vô cùng trang trọng, thành kính và\dots\ có chút gì đó thi vị quá mức.

Tôi lúng túng gãi đầu, quay sang hỏi nhẹ: ``\vie{Cái đó\dots\ E-Em vừa gọi anh là `Onii-sama' á? Nghe hơi quá rồi đó. Anh không nghĩ mình đáng được gọi như thế đâu.}''

Cô bé bối rối cười, nhưng vẫn lắc đầu đầy quả quyết: ``\vie{Không đâu ạ! Với em, Onii-sama là người tốt nhất, là người em tin tưởng tuyệt đối. Em muốn gọi như thế, vì em thấy mình có anh trai thật sự rồi.}''

Một câu nói thôi mà khiến tôi nghẹn cổ họng. Tôi mỉm cười, lòng chợt thấy ấm lại kỳ lạ. Không còn gì để phản bác nữa, tôi khẽ gật đầu: ``\vie{Vậy thì, được thôi. từ giờ anh sẽ cố gắng không làm em thất vọng, Suzuki-chan.}''

``\vie{Suzuki.}''

``\vie{Hả?}''

``\vie{Em muốn được gọi như thế. Không có `-chan' đằng sau nữa, nghe vẫn xa cách lắm ạ,}'' cô bé cau mày, phúng phính đôi má lên.

Tôi thở dài một hơi, chấp nhận luôn cách gọi đó. ``\vie{Được rồi, Suzuki.}'' Tôi quay sang Kaigou, hỏi tương tự: ``\vie{Tiện thể đang hỏi về xưng hô, cô muốn tôi xưng như thế nào?}''

Cô nàng lảng đi, đáp cụt lủn: ``\vie{\dots Gì cũng được.}''

Tôi bật cười, bông đùa một câu: ``\vie{Vậy thì tôi sẽ gọi cô là Kaigou-chan nhé?}''

``\vie{C-Cái gì chứ!?}'' cô nàng gần như bật nảy lên, mặt đỏ như cà chua, lắp bắp, ``\vie{Đ-Đ-Đ-Đợi đã, cậu nghiêm túc đấy chứ!?}''

Tôi cười khẽ, nhún vai: ``\vie{Hai đứa mình bằng tuổi nhau mà. Không được sao?}''

Kaigou trừng mắt nhìn tôi, rồi quay phắt đi, hai má vẫn hồng rực. Tôi nghĩ cô ấy không quen bị gọi như thế, nhưng mà lần đầu mà, ắt rồi sẽ quen thôi.

Sau một hồi chần chừ, cô nàng bật ra câu: ``\vie{\dots Thôi được. Nhưng đổi lại, tôi sẽ gọi cậu là Kuon-kun, nghe chưa?}''

``\vie{Rất sẵn lòng, Kaigou-chan\~{}} tôi kéo dài giọng trêu chọc, khiến cô nàng gần như muốn lăn ra đất vì ngượng, đầu bốc khói như tàu hỏa. Trời đất, trêu cô nàng có nắm đấm thép kia càng làm tôi muốn chọc thêm, như thể trêu chó dữ vậy.

\pov{Hikawa Kaigou}

Tôi vẫn chưa quen được với cái cách mà cậu ta gọi tôi là ``Kaigou-chan.'' Kể cả có là bằng tuổi đi chăng nữa, cái hậu tố đó\dots\ tôi chưa từng nghe ai dùng với mình cả. Những gã đàn ông mà tôi từng biết chỉ gọi trống không, hoặc đầy vẻ khinh khỉnh, xoi mói kiểu ``người đẹp ơi'' hay ``cô em'', nghe mà ngứa cả tai. Còn cậu ta, lại xưng hô với tôi như thể tôi là một cô bạn gái thân thiết nào đó. Vừa khó chịu, mà vừa kỳ lạ thay, cũng khiến tôi thấy vui.

Cảm giác được ai đó gọi mình một cách thân tình, nhẹ nhàng, không có ẩn ý, không có định kiến, không gợi lại bóng tối, là lần đầu tiên tôi có được.

Tôi liếc sang cô em gái của mình, và rồi một ý nghĩ hiện ra trong đầu. Trước khi Kuon kịp mở miệng thêm điều gì nữa, tôi nắm tay em gái và quay người bước đi thật nhanh. ``\vie{Đi thôi, Suzu. Kệ cậu ta đi.}''

``\vie{N-Này! Hai người định bỏ tôi lại hả!?}'' tôi nghe giọng cậu ta gọi từ phía sau.

Suzu quay lại cười tinh nghịch, còn tôi thì giả vờ không nghe thấy, kéo con bé chạy chậm dần thành bước đi đều.

Lúc đó, tôi nghe thấy tiếng Kuon cười khẽ phía sau: ``\vie{Trả thù ngọt thật đấy, Kaigou-chan.}''

Tôi không quay đầu lại. Nhưng, một nụ cười mỉm tinh nghịch nở trên môi tôi. Một nụ cười thực sự.

Tôi mong chờ lời hứa từ cậu đấy, Hikawa Kuon-kun! Đừng để làm cho tôi thất vọng!

\ct

\pov{Hikawa Kuon}

Chúng tôi tiếp tục cuộc hành trình trong The Void, ba người lặng lẽ bước giữa một không gian dường như bất tận và không có phương hướng. Suzuki ngồi yên trên lưng tôi, hai tay vòng qua vai tôi, mái đầu nhỏ tựa nhẹ vào lưng, thỉnh thoảng cựa mình như để đổi tư thế. Còn tôi và Kaigou-chan thì vừa đi vừa trò chuyện.

Chúng tôi nói đủ thứ chuyện trên trời dưới đất, từ những thói quen thường ngày đến sở thích riêng tư. Càng nói, tôi càng thấy bất ngờ vì cô ấy có nhiều điểm giống tôi đến kỳ lạ: cô ấy cũng thích máy tính, mê âm nhạc, có thể học tập và làm việc một cách nghiêm túc nếu muốn, và nói tiếng Anh lẫn tiếng Nhật đều rất tốt, thế nên cô nàng mới hiểu ý tôi nói khi lựa chọn giữa việc quay lại thế giới cũ và về với tôi.

Kaigou-chan nói rằng việc nói được hai thứ ngôn ngữ đó là do ``bẩm sinh'', nhưng mà tôi thấy khá khó tin. Cơ mà, việc cô ấy tự học trong khi phải làm việc tới mệt người để nuôi em bé cũng là một điều tôi thấy khá bất khả thi. Nhưng thôi kệ, quan tâm tiểu tiết làm gì nhỉ?

Nói thật, gọi cô ấy là ``bản sao'' của tôi cũng chẳng sai là bao.

``\vie{Vậy nếu cô muốn lập trình thì sẽ làm gì?}'' tôi hỏi trong lúc mắt vẫn nhìn về phía trước.

Kaigou-chan im lặng vài giây trước khi trả lời, như thể đang nhẩm tính điều gì đó trong đầu. ``\vie{Thường thì tôi mượn máy của mấy cô bà hàng xóm thôi, chúng cũng là mấy loại cũ mèm áy mà,}'' cô trả lời, giọng bình thản nhưng thấp đi một chút. ``\vie{Tôi vẫn chưa có đủ tiền để mua một cái đâu, bao nhiêu tiền tôi đều đổ dồn hết vào Suzu rồi còn gì.}''

Cô nàng hướng mặt về tôi, vẻ mặt thoáng trầm buồn: ``\vie{Mà nếu có viết được code thì cũng chỉ làm được mấy cái đơn giản như `Hello World' hay vòng lặp thôi, không đến mức cao siêu như cậu được.}''

Tôi khẽ gật đầu, cố không để lộ vẻ ái ngại. ``\vie{Ra thế.}''

``\vie{Mà, đời là thế,} cô cười nhẹ, dù có vẻ buồn nhiều hơn vui.

``\vie{Cả hai người nghe như hacker á\~{}}'' Suzuki bỗng lên tiếng từ trên lưng tôi, giọng pha lẫn ngạc nhiên và ngưỡng mộ. ``\vie{Em mà biết dùng máy tính giỏi như anh chị thì chắc đời khác lắm rồi.}

Tôi bật cười. ``\vie{Em mà học thêm vài năm nữa là vượt luôn Onee-sama của em ấy chứ đùa.}

``\vie{Cái đó\dots\ chắc em phải cố gắng học mười năm nữa quá.}'' cô bé nói lí nhí, rồi dụi đầu nhẹ vào lưng tôi như thể đang tìm một chút ấm áp giữa cái lạnh lẽo của không gian xung quanh.

Tôi khẽ nhìn sang Kaigou-chan. Ánh mắt cô hướng về phía trước, nhưng tôi thấy khóe môi cô hơi nhếch lên. Dù chỉ là một nụ cười thoáng qua, nhưng tôi biết, nó thật lòng.

Khi cả ba tiếp tục dạo bước giữa không gian tối đen tĩnh mịch ấy, Suzuki vẫn ngoan ngoãn ngồi trên lưng tôi, thỉnh thoảng nghịch nhẹ sợi dây đeo đồng hồ đeo tay của tôi như một đứa trẻ rảnh rỗi tìm thứ gì đó để bận tay. Kaigou-chan đi bên cạnh, tay vẫn đút túi váy, ánh mắt đảo qua lại như đang cân nhắc một câu hỏi.

``\vie{Này,} cô chợt lên tiếng, ngập ngừng, ``\vie{cậu từng nói mình là tổng quản lý gì đó ở Nhật, đúng không? Kiêm luôn vài việc lặt vặt?}''

Tôi gật đầu. ``\vie{Ừ. Cô tò mò à?}

``\vie{Tò mò thì sao nào? Bộ tôi không được hỏi chắc?} cô nàng chống hông, nhíu mày đáp lại.

Tôi phì cười. ``\vie{Tất nhiên là có chứ. Vậy\dots\ cô có biết VTuber không? Kiểu YouTuber ảo ấy?}''

Kaigou-chan khựng lại nửa bước rồi gật đầu khẽ. ``\vie{Tôi có nghe qua. Mấy nhóm lớn thì phải\dots\ \hll\ với Nijisanji gì đó, đúng không?}''

``\vie{Pín pon, đúng bài luôn,}'' tôi chỉ tay lên như thể xác nhận đáp án đúng trong trò chơi truyền hình. ``\vie{Tôi làm việc cho bên Tam Giác Xanh đấy.}''

``\vie{\dots Cái gì cơ?}'' cô nàng bật cười, như thể không tin lời tôi nói. ``\vie{Cậu mà làm cho \hll\ á?}''

``\vie{Nghe thì khó tin nhưng đúng là thế,}'' tôi nhún vai. ``\vie{Mỗi tội không phải theo kiểu tôi mong đợi. Ban đầu tôi chỉ định vào làm mảng máy tính thôi, nhưng không biết ăn nhầm bùa mê thuốc lú gì mà lại bị nhét luôn vào vị trí tổng quản lý.}''

``\vie{Số cậu đỏ thế còn gì!}'' Kaigou-chan vừa nói vừa phá lên cười, vẻ không tin nổi. ``\vie{Người ta xin được chân IT còn không xong, cậu lại vớ luôn cả chức to đùng!}''

``\vie{Cũng chẳng buồn cười gì cho cam đâu,}'' tôi lẩm bẩm, cười chát.

``\vie{Em cũng tưởng anh làm kiểu gì đó\dots\ bình thường cơ,}'' Suzuki chen vào, giọng vẫn pha lẫn ngây thơ, ánh mắt toát lên vẻ ngưỡng mộ. ``\vie{Ai ngờ lại xịn thế\dots}''

``\vie{Tôi không nghĩ là một người làm việc cật lực suốt ngày như Kaigou-chan lại biết mấy cái đó đấy,} tôi nhìn sang Kaigou, nửa đùa nửa thật.

Cô nàng phồng má phụng phịu. ``\vie{Ít nhất tôi cũng có thời gian nghỉ ngơi chứ! Cậu tưởng tôi là robot chắc?}''

``\vie{Thế mới ngạc nhiên chứ.}''

``\vie{Nhưng mà, Onii-sama, công việc của anh cụ thể là làm gì thế ạ?}'' Suzuki hỏi nhỏ, có vẻ tò mò thật sự. ``\vie{Em chỉ mới biết VTuber qua mấy đoạn clip ngắn thôi.}

``\vie{Tôi cũng muốn biết đó, Kuon-kun,}'' cô chị gái ở bên cạnh bất chợt lên tiếng, nét mặt hơi hiếu kỳ.

``\vie{Nói một cách đễ hiểu là thế này\dots}'' tôi thở dài một hơi, rồi bắt đầu liệt kê ra những phần mà tôi hay làm. ``\vie{Tôi là quản lý toàn bộ các mảng, từ hậu cần, kỹ thuật, phân phối thiết bị, quản lý tài chính, đôi khi còn phải làm cả lịch trình và hỗ trợ nội dung nữa. Nhưng thứ đáng kể nhất\dots\ chắc là mỗi Chủ Nhật tôi đều phải làm `người trông trẻ bất đắc dĩ' vì các thành viên nhà đó phá quá. Được cái về sớm.}''

``\vie{Nhiều đến thế á!? M-Mà, cái trông trẻ kia cậu nghiêm túc đấy chứ!?}'' Kaigou-chan chớp mắt.

``\vie{Tôi không nói đùa,}'' tôi cười một cách ngán ngẩm. ``\vie{Họ đúng là idol đấy, nhưng cũng là những nàng hề chính hiệu. Mỗi lần họ bày trò gì đó kỳ quặc đến mức mọi thứ vượt tầm kiểm soát, chính tôi với đồng nghiệp phải đứng ra giải quyết chứ ai.}''

Cô chị cười ngặt nghẽo: ``\vie{Mỗi tuần họ cứ phá đến thế, chắc họ nghịch ngợm lắm nhỉ?}''

``\vie{\dots Chuẩn không cần phải chỉnh,}'' tôi giương mắt lên trên cao, chẹp miệng một tiếng đầy bất lực. Nhất là tuần vừa rồi thôi, Shion-chan đã tạo ra một thứ gọi là ``miếng dán kỳ diệu,'' chỉ để mọi thứ nổ tung, gián tiếp kéo tôi vào nơi này.

Phải mất một lúc, Kaigou-chan mới bình tĩnh trở lại sau tràng cười đã đời ấy. Cô nàng tiếp tục hỏi, miệng vẫn nở nụ cười: ``Thế họ nghịch như vậy, cậu cũng than phiền như thế\dots\ nhưng cậu có thích công việc đó không?''

Tôi đứng khựng lại nửa bước. Câu hỏi tưởng đơn giản, nhưng lại khiến tôi mất vài giây ngẫm nghĩ. Tôi ngập ngừng một chút rồi mỉm cười.

``\eng{\dots Có chứ.}''

Không khí chậm lại trong một khoảnh khắc ngắn, như thể tôi đang được dịp nhìn lại hai năm làm việc vất vả nhưng đầy màu sắc của mình. Những đêm không ngủ vì lỗi hệ thống, những lần gào lên khi cả văn phòng bị đổ sập, những pha ``cầu trời'' vì lịch trình dồn dập, và những tiếng cười giòn tan khi nhìn các em ấy phá phách mà tôi chẳng thể làm gì được ngoài đầu hàng\dots

Tôi không nói ra, nhưng tôi biết: tôi chưa từng hối hận khi bước vào nơi ấy. Và có lẽ, đó là điều khiến tôi trở thành tôi của hôm nay.

Tôi vẫn đang mải nghĩ về những ngày tháng làm việc ở Tam Giác Xanh, về những trò hề của các em ấy mà giờ nghĩ lại vẫn buồn cười đến phát điên.

``\vie{Vậy chắc các gái ở đó thích cậu lắm nhỉ?} Kaigou-chan đột ngột hỏi, giọng nhẹ như gió thoảng.

Tôi khựng lại nửa bước. Câu hỏi nghe thì có vẻ đơn giản, nhưng lại làm tôi thấy hơi lúng túng.

``\vie{Thú thật thì\dots\ tôi cũng chẳng biết nữa,} tôi gãi má, nhìn sang hướng khác như thể tránh ánh mắt cô ấy.

``\vie{Hể\dots?} Suzuki nghiêng đầu sau lưng tôi, giọng ngạc nhiên. ``\vie{Em tưởng các chị ấy sẽ tranh nhau giành lấy Onii-sama cơ chứ.}''

Tôi bật cười, nhẹ thôi, nhưng là thật lòng.

Trong thâm tâm, tôi chưa bao giờ nghĩ mình là người đáng để ai đó phải tranh giành cả. Tôi chỉ là một tổng quản, một nhân viên kỹ thuật lặng thầm làm việc sau cánh gà. Những lần tôi viết kịch bản, lập bảng thanh toán, chỉnh dây micro, kiểm tra ánh sáng, canh giờ phát sóng trực tiếp, hô to: ``Năm giây nữa bắt đầu!'', tất cả đều để họ có thể tỏa sáng trên sân khấu. Tôi là người giữ ổn định mọi thứ phía sau, để các cô gái đó có thể thỏa sức thể hiện cá tính bản thân trước mắt hàng nghìn, hàng triệu người hâm mộ.

Tôi nhớ lại sinh nhật mình bốn tháng trước\footnote{Xem lại Quyển Birthday.}. Các cô gái bất ngờ tổ chức tiệc, còn giấu tôi kỹ đến mức không hé ra một lời. Họ hát, họ cười, và tặng tôi món quà nhỏ mà tôi thầm mong, hoặc là do chính tự tay họ tạo nên. Họ nói: ``Không có anh thì bọn em không thể yên tâm mà livestream được đâu,'' hay ``Kuon-senpai cũng là một phần quan trọng của team đó!''

Nhưng tôi không chắc mình tin hẳn vào điều đó.

Tôi biết, đôi khi lời nói tốt đẹp được thốt ra chỉ để làm vui lòng nhau. Tôi cũng biết, ranh giới giữa một ``người được quý mến'' và một ``người đồng hành thực sự'' là rất mong manh. Tôi thấy mình có lẽ chỉ là một công cụ tốt. Có ích. Nhưng chưa đủ quan trọng để thuộc về sân khấu ấy. Thứ ánh sáng mà họ mang theo, nó vượt xa tôi. Một thứ có thể dễ dàng thay thế nếu công cụ đó quá già cỗi và không còn hiệu quả.

Tôi khiêm tốn đến độ có lẽ là tự phụ. Và tôi biết điều đó.

``\vie{Vậy à\dots}'' Kaigou-chan lên tiếng, kéo tôi khỏi dòng suy nghĩ. ``\vie{Thế những người trong công ty nói như thế nào với cậu?}''

``\vie{Thú thật, tôi được A-san - đồng cấp của tôi - với Tanigo-san nói rằng tôi được các thành viên tài năng yêu quý lắm, nhưng mà\dots}

``\vie{Để tôi đoán,} cô nàng nheo mắt, nhếch mép. ``\vie{Cậu không để ý, đúng không?}

Tôi cười gượng, cố bao biện ra một lý do: ``\vie{Ham việc quá nên mới như vậy. Kiếm tiền nuôi gia đình là chính mà. Sao cô lại hỏi câu đó?}''

Kaigou nhún vai, hai tay đan ra sau, tỏ ra như thể đó chỉ là một câu hỏi vu vơ. ``\vie{Trực giác con gái nhạy bén lắm. Tôi chỉ thấy cậu là kiểu người\dots\ nói sao nhỉ, không nhận ra được ai đang nhìn mình bằng ánh mắt đặc biệt thôi.}

``\vie{Em cũng thấy thế,}'' Suzuki ngáp dài một tiếng. ``\vie{Mà nếu các chị ấy không thấy anh đặc biệt\dots\ thì em thấy đó, Onii-sama.}''

Tôi phì cười, lắc đầu bất lực.

``\textit{Hai người đúng là\dots}'' tôi nghĩ, rồi chợt thấy lòng mình dịu lại. Những lời ấy, dù đơn sơ, lại khiến tôi ấm áp lạ kỳ. Có lẽ\dots\ mình không vô hình như mình tưởng.

\pov{Hikawa Kaigou}

Tôi liếc nhìn cậu con trai đang bước bên cạnh mình, mái tóc đen nhánh xù xì rủ xuống trán, ánh mắt như đang thả trôi vào một dòng suy nghĩ nào đó xa xôi. Cái cách mà Kuon-kun lúng túng khi tôi hỏi chuyện ``các gái có thích cậu không'' khiến tôi bật cười nhẹ. Nhìn kiểu gì thì cũng rõ là cậu ta hoàn toàn không giỏi trong mấy chuyện tình cảm.

Một đứa con trai có thể thấu hiểu được tâm lý của tôi, nắm bắt được cả những khao khát mà tôi cố giấu suốt bao năm, lại trở nên vụng về khi nói đến cảm xúc của những cô gái khác. Thật kỳ lạ.

Không phải cậu ta vô tâm đâu. Tôi biết. Chẳng qua là cậu đặt nặng cái gọi là trách nhiệm lên vai mình quá, nên mới cố tình lờ đi mấy chuyện mà cậu cho là ``vặt vãnh'' như tâm sinh lý hay cảm xúc con gái thôi. Mà nghĩ lại, cách mà cậu ta thu phục tôi\dots\ có phải là vì trách nhiệm không nhỉ? Cũng có thể lắm chứ.

Nhưng đó cũng là điều khiến cậu ta khác biệt hẳn, và cũng là điều khiến cho tôi thích thú.

Tôi ngẫm một lát, rồi buột miệng: ``\vie{Mà tôi khuyên thật, cậu cũng nên để ý họ nhiều hơn nữa đi.}''

Kuon-kun ngoảnh sang, rướn mày: ``\vie{Sao không để ý? Gần như lúc nào tôi chẳng dõi theo họ? Thậm chí là ngồi ở nhà tôi cũng nghĩ tới họ chứ.}''

Tôi buông một tiếng thở dài. Rõ ràng cậu ta đần hơn tôi tưởng. ``\vie{Vấn đề là, cậu chỉ để ý tới công việc của họ thôi. Tâm sinh lý của họ thì sao? Tôi có cảm giác cậu chỉ làm vì nó là trách nhiệm của mình ấy.}''

``\vie{Thì bởi vì trách nhiệm của quản lý là vậy mà,}'' cậu ta nhún vai, giọng đều đều như một lời biện hộ.

Tôi bĩu môi, tặc lưỡi một cái rõ to: ``\vie{Sai bét! Cậu không chỉ là người giữ nhịp công ty. Cậu phải là một quản lý vừa có tầm, vừa có tâm nữa!}''

``\vie{Ý cô là sao?}'' Kuon-kun dừng lại, đôi mắt ánh lên vẻ nghi ngờ.

Tôi xoay người đi lùi mấy bước, đối diện với cậu ta, rồi giang tay ra như đang giảng đạo lý: ``\vie{Các cô gái ở đó cần cảm giác được quan tâm, được chăm sóc. Không phải kiểu giám sát hay bắt lỗi thông thường đâu. Mà là từ trái tim cơ. Cái sự giúp đỡ ấy phải đến từ sự đồng cảm, chứ không phải bảng phân công nhiệm vụ.}''

``\vie{Tôi có làm vậy mà. Mỗi khi họ có gặp điều gì khó khăn, họ luôn tìm tới tôi đấy chứ,}'' cậu ta phản bác nhẹ nhàng.

``\vie{Có, nhưng chưa đủ.}'' Tôi nhấn mạnh, giọng hơi cao lên như người thúc giục: ``\vie{\textbf{Cậu phải làm nhiều hơn như thế!}}''

Kuon-kun thở dài, nở một nụ cười chát mà tôi không chắc có phải là thương hại bản thân không. ``\vie{Cô nói như thể cô là nhân viên làm ở đó vậy\dots}''

Tôi bật cười: ``\vie{Ờ thì, con gái với nhau mà. Tôi hiểu cảm giác đó. Tôi biết điều này nghe có vẻ hơi ảo ma, nhưng nghe tôi nói này,}'' tôi ngừng lại một nhịp, hạ giọng xuống, ``\vie{con gái vốn dĩ có thể rung động từ cái nhìn đầu tiên đấy.}''

Cậu ta nheo mắt lại, rồi đáp lại một cách miễn cưỡng: ``\vie{\dots Nghe xàm xí thế nào ấy. Đời chứ có phải là anime hay manga quái đâu mà.}''

``\vie{\textbf{Thế mà có thật đấy!}}'' tôi cười tươi rói. ``\vie{Nhưng điều đó có nghĩa là: khi một cô gái yêu một chàng trai, họ sẽ luôn tin tưởng rằng người đó đang làm mọi thứ vì mình. Vì vậy, để đáp lại lòng tin đó, con trai phải thể hiện sự quan tâm của mình bằng hành động, bằng ánh mắt, bằng sự lắng nghe\dots\ không chỉ là công việc!}''

``\vie{Tôi đang nghe đây,} Kuon-kun đáp, giờ thì cậu ta đã thôi cãi. ``\vie{Tôi phải làm gì?}''

``\vie{Đơn giản lắm. Chỉ cần quan tâm tới cảm xúc họ nhiều hơn. Chứng minh rằng tấm lòng của họ không bị đặt nhầm chỗ.}''

``\vie{Ý cô là phải nắm bắt tâm lý á?}''

``\vie{Pín-pon! Chuẩn không cần chỉnh!}'' tôi nhại lại đúng theo cách mà cậu ta vừa nói.

``\vie{Sao tôi thấy khó thế nhỉ\dots}'' cậu ta gãi đầu, miệng nửa cười nửa méo.

``\vie{Cứ làm việc với họ nhiều hơn nữa, càng gần họ càng tốt}'' tôi nheo mắt trêu, đưa ra một lời khuyên mà tôi nghĩ cậu ta sẽ hiểu. ``\vie{Cậu sẽ dần hiểu rõ họ thôi. Giống như cách mà cậu đã làm với tôi và Suzu ấy. Nếu không thấu hiểu và kiên nhẫn, cậu đâu có thể thuyết phục được bọn tôi, đúng không?}''

``\vie{Nhưng kể cả cô nói thế thì-}'' Kuon-kun định lên tiếng lập luận, nhưng rồi Suzu đột ngột xen vào từ sau lưng cậu: ``\vie{Em đồng tình với Onee-sama! Onii-sama đúng là có siêu năng lực đọc cảm xúc đấy! Nhưng mà\dots\ chỉ khi anh ấy thực sự để ý thôi!}''

``\vie{Đến Suzu cũng đồng tình với tôi đấy thôi!}'' tôi khoanh tay, mặt vênh lên. ``\vie{Đừng tưởng cậu giỏi IT rồi cái gì cũng có thể giải quyết bằng logic hết nhé.}''

Kuon-kun thở dài, nhưng nụ cười nhẹ vẫn lộ ra nơi khóe môi.

Tôi biết cậu ta không hẳn đã bị thuyết phục hoàn toàn, nhưng chí ít, tôi đã gieo vào đầu cậu một mầm mống của sự thay đổi. Nếu cậu có thể làm được với hai chị em chúng tôi, cậu cũng có thể làm thế với những người làm việc với cậu đó, Kuon-kun.

Tôi ngẩng mặt lên nhìn bầu trời tím sẫm lơ lửng bên trên, thứ trần nhà bất định của cái nơi hư vô kỳ quái này, rồi khẽ nghiêng đầu nhìn sang cậu con trai.

Cậu ta vẫn đi song hành bên tôi và Suzu, tay thọc túi áo, bước chân đều đều như một con lắc đã lên dây cót. Ánh mắt cậu ta không còn đượm vẻ trầm tư như ban nãy nữa, nhưng tôi vẫn nhận ra có điều gì đó chưa nói ra.

Tôi giả vờ như hỏi vu vơ: ``\vie{Mà này, Kuon-kun\dots\ cậu tới đây bằng cách nào vậy?}''

Cậu ta hơi sững lại, rồi thở dài một tiếng: ``\vie{Câu hỏi hay đấy. Tôi cũng đang tự hỏi điều tương tự từ lúc rơi xuống đây.}''

``\vie{Là sao?}''

``\vie{Cũng giống như hai chị em cô thôi. Bị một cái vòng xoáy bí ẩn bất chợt xuất hiện, hút tôi vào trong, đưa tôi tới một nơi của nợ này.}''

``\vie{Ý tôi không phải thế,}'' tôi cau mày, nhìn thẳng vào đôi mắt cậu ta. ``\vie{Từ lúc mọi chuyện dẫn tới cái lỗ xoáy ấy.}''

Kuon-kun nắm chặt hai chân của Suzu, rồi nhắm mắt lại, bắt đầu chậm rãi kể lại mọi thứ: ``\vie{Chuyện là sáng nay, tôi tới văn phòng làm việc để đưa cho A-san một tờ đơn xin nghỉ việc, nhưng đó không phải của tôi, mà là tờ chị ấy nhờ tôi inn ra để đưa hộ cho người khác thôi.}

``\vie{Sau đó thì sao?}''

``\vie{Hầy\dots\ Đang định làm việc thì ông bạn tôi ở trường bảo là giáo viên bộ môn đẩy lịch bảo vệ lên hôm nay, thế là tôi ba chân bốn cẳng chạy lên trường. Nhưng mà trước đó ấy, tôi có để lại một mảnh giấy nhắn rồi.}''

Cậu ta ngừng một chút, rồi thở ra: ``\vie{Tôi viết vội. Nét chữ nguệch ngoạc. Nội dung là: ``Hôm nay anh không thể đến để gặp mặt các em được. Anh đã quá mệt mỏi với chuyện này rồi. Đừng gọi cho anh.'' đại loại thế.}''

Tôi khựng lại một nhịp. Suzu từ phía sau lưng cậu ta cũng rướn người. Cả hai chị em tôi nhìn nhau, rồi cùng đưa ra một kết luận.

\begin{dialogue}
    \speak{Suzuki} Cái đó nghe như\dots
    \speak{Kaigou} \dots một bức thư tuyệt mệnh ấy.
\end{dialogue}

Kuon-kun bật cười gượng, đưa tay gãi đầu: ``\vie{Tôi không nghĩ nó sẽ tới mức thế đâu. Chỉ là tôi không muốn bị làm phiền trong lúc đang căng thẳng với đồ án thôi.}''

Câu trả lời bâng quơ của cậu ta càng khiến tôi sốc hơn. ``\vie{Trời ơi! Cậu mà để lại mấy dòng như vậy ở công ty thì thể nào họ cũng tá hỏa lên cho xem. Cậu có biết là nó nghe như kiểu một người từ chức rồi biến mất khỏi cuộc đời không?}''

Cậu ta ngoảnh mặt về phía tôi, nheo mắt lại, nửa ngạc nhiên, nửa nghi ngờ vì cho rằng tôi đang quá lố: ``\vie{Ờm\dots\ Cô nói sao cơ?}''

Suzu nhỏ giọng chen vào: ``\vie{Nếu là em\dots\ chắc em sẽ khóc đấy.}''

Tôi gật đầu, giọng trở nên nghiêm túc hơn: ``\vie{Tôi có cảm tưởng rằng\dots\ họ đang lo lắng cho cậu đấy. Rất nhiều là đằng khác. Hành văn kiểu thế thì chết!}''

``\vie{A-nô-ne (Em thì thấy)\dots}'' cô em gái tôi chị mình, vòng tay ôm cổ Kuon-kun từ phía sau như muốn khẳng định, ``\vie{Họ quan tâm tới anh lắm, đến mức nếu không thấy anh trả lời thì họ sẽ tìm anh bằng mọi giá luôn ấy ạ.}''

Tôi nhấn mạnh phụ họa, nêu rõ tính nghiêm trọng của vấn đề: ``\vie{Với cách mà cậu tự xem mình chỉ như một người `hậu trường thầm lặng', có khi chính cậu mới là người mà họ xem là trụ cột. Vậy mà cậu để lại giấy nhắn đó mà không để ý tới tấm lòng sẽ ra sao ư?}''

Kuon im lặng không đáp. Tôi liếc sang thì thấy cậu ta nhìn xuống đất, gương mặt hơi đờ ra. Có lẽ lúc đó, cậu mới nhận ra hậu quả từ cái lá thư ngắn ngủi viết trong phút ngẫu hứng ấy.

Dẫu sao thì tôi cũng không trách. Vì tôi biết cái cảm giác bị đè nặng bởi trách nhiệm, rồi chỉ muốn biến mất tạm thời để tìm lại chính mình. Tôi cũng từng như thế. Mà, biết đâu\dots\ chính vì vậy mà tôi lại đồng cảm được với cậu ta đến thế này.

\begin{center}
    \uline{Thành phố điện tử - Vũ trụ 001}
    \pov{-}
\end{center}

Cơn lốc xoáy kia biến mất nhanh như khi nó xuất hiện.

Trung tâm Akiba ảo nay chỉ còn là một đống đổ nát lặng ngắt. Những toà nhà kính cao tầng trơ khung thép, kính tường vỡ vụn như bị một bàn tay khổng lồ cào nát. Các cột đèn loạng choạng, dây điện rối như tổ quạ. Cây cầu bắc ngang con kênh bên dưới đã gãy làm đôi, trụ chống méo mó, lưng cầu chìm trong làn khói đen vẫn chưa tan hết.

Cảm tính của hai chị em nhà Hikawa ở vũ trụ 717 đã đúng.

``Biệt đội tìm và bắt Kuon-senpai'', cái tên mà nàng Sói đen bịa ra trong lúc điên tiết vì bị cậu Tổng bỏ đi không lời từ biệt, lúc này đang len lỏi qua những con đường tan hoang. Cả nhóm năm người tiến dần về phía cây cầu vừa bị quét qua bởi lỗ đen.

``Không\dots\ thể nào\dots'' Towa thốt lên, mắt mở tròn đầy hoảng hốt.

``Đùa nhau à?'' Mio lẩm bẩm, giọng không giấu nổi hoang mang.

``Tại sao chứ?'' Watame siết chặt hai tay trước ngực, gương mặt tái nhợt.

Trước mắt họ là một cảnh tượng tang thương: kính vỡ ngổn ngang, xe cộ bị hút dạt ra, một vài cột điện còn phát ra tiếng nổ lách tách. Không khí vẫn nồng mùi khét của khói cháy.

Nàng Ác ma rút điện thoại ra, tra định vị chiếc xe buýt mà Kuon từng nói sẽ dùng để đi từ công ty về. Dấu chấm đỏ hiện lên, đứng yên giữa một khúc cầu đã đổ sập.

Cây cầu đó - nơi mà họ đang đứng trước mắt - chính là nơi bị ảnh hưởng nặng nề nhất do cái lỗ đen kia.

``Lẽ nào\dots\ Kuon-senpai đã\dots!?'' Noel thốt lên, môi run rẩy.

``Bình tĩnh nào, Noel,'' Mio đặt tay lên vai cô, liền trấn an, ``anh ấy sẽ ổn thôi.''

``Mọi người, mau đi thôi!'' Haato hối thúc nhẹ, giọng run nhưng vẫn quyết đoán.

Cả nhóm gật đầu, chạy dọc theo tuyến đường bị phong tỏa. Họ len lỏi qua đống đổ nát, tìm tới chiếc xe buýt đã lật nghiêng. Đám đông hiếu kỳ, cảnh sát, lính cứu hộ chen chúc nhau tạo thành một bức tường người.

``Xin lỗi, cho chúng tôi đi qua\dots'' họ vừa luôn miệng xin phép, vừa đẩy nhẹ cắt ngang để tiếp cận tới chiếc xe.

Giữa lúc đó, một gã đàn ông mập ú tranh thủ trong đám người hỗn loạn kia đã nhéo mông Haachama. Cô nàng rít lên: ``\textbf{Dám sàm sỡ ta hả, cái tên lợn mập này!?}''

``Kiềm chế lại đi, Haato-san! Họ không cố ý đâu!'' nàng Sói đen hốt hoảng can ngăn, giữ lại hai vai của người tiền bối.

``\textbf{Cố ý cái đầu họ! Bỏ ra, chụy xử lý-}'' cô nàng Song tính xắn tay áo lên, định đấm cho gã ta một cái, nhưng bị nàng Ác ma níu cô lại, giật mạnh áo: ``Giờ không phải lúc đâu chị ơi!''

Sau một hồi kìm chế, họ cũng tới được chiếc xe buýt. Xe cứu thương chật kín, lính cứu hỏa bắc thang, dỡ kính, đưa các nạn nhân ra, cả người sống lẫn những người không còn thở.

Mọi người đổ dồn ánh nhìn vào từng người được khiêng ra. Không có Kuon, hoặc ít nhất là có ai đó mang bóng dáng của cậu.

Trong lúc cả nhóm tiếp cận gần hơn tới chiếc xe buýt còn đang bị lấp xuống đống ngổn ngang của đá và bê tông, nàng Cừu bất chợt vấp ngã. ``Ái cha!''

``Bà không sao chứ, Watame!?'' Towa hỏi vội.

``T-Tôi ổn\dots'' nàng Cừu gật đầu, nhưng tay vẫn chống đất, mặt lấm lem bụi đen.

Nữ Hiệp sĩ bất chợt cất tiếng run rẩy: ``M-Mọi người ơi\dots'' Ánh mắt mọi người đổ dồn về phía mà nàng Kỵ sĩ đang chỉ.

Từ một khung cửa kính vỡ nát của xe buýt, một cánh tay thò ra, bất động. Trên cổ tay ấy là một chiếc đồng hồ thông minh mặt tròn, màu thép đen bạc.

Chiếc đồng hồ đó\dots\ giống hệt cái mà cậu Tổng quản lý Kuon thường đeo mỗi ngày đi làm.

Cả nhóm sững sờ.

``Đ-Đừng nói là\dots'' Haato thì thào, mặt không còn giọt máu.

``Anh ấy\dots'' Watame bắt đầu bật khóc.

``Chết rồi sao\dots?'' Towa nén nghẹn, nhưng giọng vỡ ra ở cuối câu.

``K-Không thể nào\dots'' Mio cắn môi, mắt đỏ hoe.

``Chúng ta đã làm gì sai\dots?'' Watame nấc lên.

``Tại sao\dots\ \textbf{Tại sao anh lại phải chết chứ!?}'' Noel thét lên, nước mắt tuôn như mưa.

Thời gian như đông cứng lại, tiếng còi cứu thương, tiếng người hô hoán, tiếng kim loại va chạm lách cách, tất cả mờ đi, chỉ còn lại nhịp tim đập gấp gáp, nỗi bàng hoàng nặng nề và một luồng cảm xúc lạnh ngắt siết lấy lồng ngực họ.

Không khí trở nên đặc quánh. Ngột ngạt. Như có thứ gì đó vô hình đang siết chặt lấy cổ họ, không cho họ thở. Mỗi người đều mang một biểu cảm khác nhau, nhưng đều chung một nỗi đau khôn tả: người mà họ hết mực quý trọng đã không còn.

Một cơn gió lướt qua, mang theo hơi lạnh cuối thu của bầu trời xám xịt phía trên.

Và rồi, những giọt mưa đầu tiên bắt đầu rơi.

Tí tách. Tí tách.

Nhẹ, rồi nhanh, rồi nặng hạt dần. Mưa như trút, ầm ầm xối xuống mặt cầu đã nứt toác, hòa cùng nước mắt đang rơi lã chã từ những cô gái vừa mới ngày hôm qua còn cười đùa, nay đã rũ rượi như những con búp bê bị bỏ rơi.

Nàng Cừu run lẩy bẩy, tay siết chặt ngực áo. Nàng Ác ma đứng bất động, tay nắm chặt thành nắm đấm như dằn vặt bản thân. Nàng Song tính vừa khóc vừa lẩm bẩm tên cậu như một câu thần chú. Nàng Sói đen đứng lặng thinh, môi mím chặt. Còn nàng Hiệp sĩ thì quỳ sụp xuống, đôi vai rung lên từng hồi như bị sét đánh.

Mưa rơi mỗi lúc một nặng. Trên cây cầu chằng chịt vết nứt, nước đọng lại từng vũng, khiến mặt đường trở nên trơn trượt. Cảnh sát bắt đầu rào lại khu vực, yêu cầu người dân rút lui khỏi vùng nguy hiểm.

Tiếng kim loại rên rỉ vang lên.

*RẮC*

Một tiếng bê tông, hoặc một tiếng kim loại bị bẻ gãy.

Một đoạn của cầu sụt xuống vài centimet. Những người lính cứu hộ lập tức hô hoán:

``\textbf{Cầu có thể sập bất cứ lúc nào! Tất cả rút lui ngay!}''

``\textbf{Lui lại! Tất cả lui lại!}''

``\textbf{Đừng đứng gần khu vực đó nữa!}''

Mọi người giật mình quay đầu, hoảng hốt, chạy về phía sau.

Nhưng trong khoảnh khắc đó, Towa và Watame bất ngờ lao về phía cái xác vẫn còn nằm kẹt giữa cửa kính, nơi có cánh tay đeo đồng hồ quen thuộc thò ra.

``\textbf{Không! Anh ấy vẫn còn ở đó!} nàng Cừu hét lên, nước mắt hòa với mưa.

``\textbf{Chúng tôi sẽ đưa anh ấy ra!}'' nàng Ác ma gào lên, cố vùng khỏi sự níu kéo.

``Không được! Lui lại! Nguy hiểm lắm!'' một người lính cứu hỏa giữ chặt vai cả hai.

``Buông ra! Anh ấy còn sống mà! Tôi xin các người\dots!''

Cầu bắt đầu rung lắc. Một phần dầm thép rơi xuống sông phía dưới, va chạm làm nước bắn tung tóe.

Mio và Noel hoảng hốt kéo hai cô em về.

``\textbf{Đừng có điên như thế! Hai người cũng muốn chết theo Kuon-senpai à!?}'' nàng Sói đen quát lên, nước mắt vẫn giàn giụa.

``\textbf{Nhưng anh ấy\dots\ anh ấy\dots!}'' Watame nức nở.

``Chúng ta không chắc đó là Kuon-senpai đâu! Làm ơn\dots\ xin đừng\dots'' nữ Hiệp sĩ vừa khóc vừa giữ chặt Towa.

Một tiếng *RẦM!* nữa vang lên.

Một phần lan can cầu sập xuống, kéo theo một mảng mặt đường. Dưới mưa, mọi thứ trở nên hỗn loạn.

Cả nhóm buộc phải lùi lại, bị kéo khỏi hiện trường trong nước mắt và tiếng gào thét. Họ bất lực nhìn lại phía xác người ấy, cái xác mà họ vẫn tin là Kuon, đang dần chìm vào làn mưa xám xịt và đống đổ nát ngày một lún sâu hơn.

Không ai có thể cầm lòng. Họ không muốn rời đi. Nhưng cây cầu thì không cho họ cơ hội chọn lựa.

\ct

Cả năm người đứng lặng bên bờ sông, mặc cho mưa xối xả trút lên đầu. Gương mặt họ tái đi vì lạnh, vì mưa, và vì nỗi đau đang gặm nhấm từ trong tim.

Họ không còn cười đùa. Không còn trêu nhau. Không còn gọi ``Kuon-senpai'' bằng giọng nũng nịu nữa.

Chỉ còn tiếng khóc hòa lẫn với tiếng mưa, và sự thật lạnh lùng rằng: người con trai đó đã không còn trên cõi đời này.

Hoặc ít nhất là, họ tin như vậy.

Bởi họ đâu thể ngờ rằng, Hikawa Kuon thật sự vẫn còn sống. Chỉ là cậu đang lạc ở một vũ trụ khác.

\begin{center}
    \uline{The Void.}
    \pov{Hikawa Kuon}
\end{center}

Mưa\dots\ Tôi không biết vì sao trong đầu mình lại vang vọng âm thanh đó. Mặc dù ở đây hoàn toàn tĩnh lặng, không có lấy một giọt nước, một cơn gió, hay cả ánh sáng mặt trời, vậy mà tôi vẫn tưởng tượng được tiếng mưa rơi lạnh buốt trên da thịt mình.

Chắc là do linh cảm. Hoặc tệ hơn, là nỗi ám ảnh mơ hồ trong tôi về thứ đang xảy ra bên ngoài, ở thế giới nơi tôi vừa rời đi.

Tôi cắn môi. Lòng thắt lại một cách khó hiểu.

Kaigou-chan vừa mới buông một câu khiến tôi vừa ái ngại, vừa lạnh sống lưng: ``\vie{Tôi thậm chí còn đang nghĩ họ đang khóc cho cậu đấy. Rất đau thương.}''

Tôi thoáng rùng mình. Dù không có bằng chứng cụ thể, tôi biết cô nàng không nói suông. Cảm giác kỳ lạ ấy đang vang vọng bên trong tôi. Một sự kết nối vô hình giữa những người thật lòng nghĩ về nhau, hoặc có lẽ, là một sự liên kết giữa những người con gái.

Tôi siết nhẹ hai chân Suzuki sau lưng. Cô bé vẫn nằm yên, thi thoảng lại tựa đầu vào vai tôi và thở khe khẽ, như thể muốn góp phần làm dịu đi cảm xúc hỗn độn trong tôi.

``\vie{Thì\dots\ cũng là vì vội mà,}'' tôi cố cười, tìm cách chống chế. ``\vie{Tôi chỉ định để lại một lời nhắn gọn gàng rồi chạy lên trường cho nhanh thôi.}''

*BỐP*

Tôi hơi choáng vì cú đập vào lưng bất ngờ. Một cái đập không đau lắm, nhưng đủ khiến tôi giật mình. Suzuki cũng ngóc đầu dậy trong giây lát.

``\vie{Cô làm cái gì đấy--}'' tôi quay lại nhìn, nhưng bắt gặp ánh mắt của Kaigou-chan đang lườm như thể tôi vừa nói điều ngu ngốc nhất thế giới.

``\vie{Cậu nói thế mà nghe được à?}'' cô chống hông, mắt nhíu lại. ``\vie{Mấy cái câu như `Anh quá mệt mỏi', `Đừng gọi cho anh nữa' mà cậu gọi là gọn gàng à? Thế thì còn thiếu mỗi câu `Tạm biệt mọi người' thôi đấy!}''

Tôi cứng họng. Không thể phản bác lại. Vì đó chính xác là tôi đã viết như thế.

``\vie{Cô\dots\ thấu hiểu con gái tới mức độ đó cơ à\dots}'' tôi lẩm bẩm.

Cô nàng khịt mũi, vênh mặt đầy tự tin: ``\vie{Không phải ai cũng có vinh dự được tôi khuyên bảo đâu nhé. Nhất là con trai. Cậu là người đầu tiên đấy.}''

Tôi nhìn cô, rồi nhìn sang Suzuki. Cô bé chỉ mỉm cười, ánh mắt dịu dàng như muốn nói: ``Onee-sama đó giờ vẫn vậy, chỉ là chưa từng có ai khiến chị chịu mở lòng thôi.''

Tôi thở dài. Tôi hiểu ra một điều.

Từ những lời mà hai chị em họ nói, từ cái cách mà Kaigou vạch ra từng lỗ hổng trong suy nghĩ của tôi, tôi bắt đầu xâu chuỗi lại tất cả những gì mình đã làm suốt hai năm qua.

Tôi từng nghĩ mình là một người quản lý giỏi. Làm việc nhanh gọn, xử lý sự cố hiệu quả, biết giữ cho toàn bộ hệ thống vận hành trơn tru.

Nhưng hóa ra, đôi khi cái tôi thiếu lại là những điều nhỏ nhặt, như việc ngồi xuống nói chuyện thật lòng, hay đơn giản là lắng nghe cảm xúc của ai đó mà không đánh giá họ qua năng suất.

Những lần tôi tư vấn vói Ru-chan, Towa-chan và Peko-chan\footnote{Xem lại {\flaregothic ÉPISODE 111}, quyển số Chín.}, hay những lần tôi an ủi Koro-chan khi bị mất chiếc máy chơi điện tử yêu thích nhất\footnote{Xem lại {\flaregothic ÉPISODE 97}, quyển số Tám.}\dots\ chừng đó vẫn chưa đủ. Tôi phải làm nhiều hơn thế.

Tôi gãi đầu, ngẩng nhìn màn sương tím mờ mịt phía trước. Có lẽ tôi nên học cách lắng nghe các cô gái nhiều hơn thật.

Tôi lâm bẩm cẩu đó, khiến Kaigou-chan dỏng tai lên. Cô nàng ngạc nhiên quay sang nhìn tôi, rồi phì cười: ``\vie{Thay đổi sớm vậy à?}''

``\vie{Tôi thích nghi tốt lắm,}'' tôi nhún vai. ``\vie{Với lại\dots\ tôi cũng không muốn một ngày nào đó, lại để mấy cô gái tôi làm việc lâu năm bật khóc trong mưa như thế.}''

Suzuki nhẹ nhàng thì thầm sau vai tôi: ``\vie{Em nghĩ\dots\ lần này Onii-sam sẽ làm tốt hơn.}''

Tôi nở một nụ cười nhẹ.

Có lẽ tôi nên như vậy thật.

\pov{Hikawa Kaigou}

Tôi không biết nữa.

Cảm giác ``dạy đời'' cho người khác, đã thế lại là con trai, là một cái gì đó khiến tôi thấy kỳ quặc.

Vừa nãy tôi nói một tràng không ngừng nghỉ như thể mình là chuyên gia tình cảm vậy. Mà đúng là như thế thật, ít nhất là trong ba phút ngắn ngủi. Còn cái tên ngốc Kuon-kun kia thì lại ngoan ngoãn nghe hết không sót một chữ, lại còn gật đầu ra chiều thấu hiểu nữa chứ.

Không biết sao, tự dưng tôi thấy vui. Cảm giác như được làm người chỉ dẫn vậy. Không phải kiểu ``chị đại dẫn dắt đám em út'', mà là một cái gì đó nhẹ nhàng hơn, dễ chịu hơn. Như thể, cuối cùng thì tôi cũng được công nhận về mặt nào đó, ngoài cái việc đấm đá và chạy trốn như một kẻ lỗ mãng trong con mắt của người thường ấy.

\il

Nhưng khổ nỗi\dots

Niềm vui ấy chẳng kéo dài được bao lâu. Vì tôi đang cảm nhận một cơn sóng ngầm\dots\ theo đúng nghĩa đen ở đó. Ngay dưới vị trí ở bụng dưới.

Chết tiệt.

Tôi vẫn còn buồn đi vệ sinh.

Tôi biết, nghe thì chẳng sang chút nào. Nhưng mà là thật! Từ lúc Suzu vô tư ``xả nước cứu thân'' hồi nãy, tôi đã bị ảnh hưởng tâm lý. Cố nín cũng nín, cố quên cũng quên, rồi bao nhiêu chuyện dồn dập làm tôi xao nhãng. Tôi tưởng là qua rồi, tưởng là tôi mạnh mẽ lắm, có thể chịu được cho tới khi tìm được một nơi để đi vệ sinh đúng nghĩa.

Ai ngờ giờ yên ổn lại, cái ``nỗi niềm'' ấy bỗng dưng trở lại, và lần này nó còn rõ rệt hơn cả hồi nãy.

Cảm giác như tôi đang gánh nguyên một cái bể chứa nước ở bụng dưới. Cứ mỗi bước đi lại là một nhịp tim tăng tốc, một sự nhắc nhở rõ ràng rằng: ``Ê, cô còn có việc chưa giải quyết kìa.''

``Híc\dots!''

Tôi cảm tưởng như tôi bắt đầu ra rồi. Vội lật chiếc váy lên, tôi nhìn thấy chiếc quần lót của mình vẫn khô cong. Thở phào.

Nhưng mà nếu cứ như thế này mãi thì cũng không được.

``\vie{\hl{Mình muốn đi tè, mình muốn đi tè, mình muốn đi tè\dots!}}'' tôi lẩm bẩm lặp đi lặp lại, hai tay cố giữ chặt lấy phần dưới, thậm chí còn tính đến chuyện\dots\ lấy hẳn tay bịt vào ``lỗ xả'' kia.

Nhưng nếu làm vậy, liệu Kuon-kun và Suzu có thể thấy không? Tôi đang đi chậm hơn một chút, ngay theo sau họ, nên chắc họ không thấy được. Cơ mà nhỡ đâu nếu họ quay đầu lại thì\dots

Ư\~{} Khó chịu quá! Mót tè quá!

Tôi lén đảo mắt nhìn quanh, mong tìm được một góc nào đó để giải quyết ``tình hình chiến sự'', nhưng mà ở cái nơi kỳ quái gọi là The Void  - theo lời của Kuon-kun - thì tìm đâu ra nhà vệ sinh? Toàn là sương mù và mặt đất phẳng lì đến phát chán. Nếu gặp may thì cũng chỉ là mấy cái vật dụng chẳng giúp được gì.

Tôi, Hikawa Kaigou - một cô gái mạnh mẽ, bền bỉ, đã sống sót qua hơn bảy năm địa ngục - giờ đây phải đối mặt với thứ thử thách tàn nhẫn nhất mà nhân loại từng biết đến:

\emph{Cơn buồn đi vệ sinh.}

Và điều tệ hơn nữa?

Tôi biết, biết chắc luôn, là có một tên nào đó ngồi sau màn hình đang vung phím lạch cạch, viết nên hoàn cảnh trớ trêu này cho tôi, bắt tôi tuôn ra những lời bẩn thỉu như thế, rồi còn nhếch mép cười như thể đây là đỉnh cao của sáng tạo văn chương.

Tôi không ngu đâu. Tôi thấy ánh mắt cậu đó - ừ, đúng cậu đấy, cái người đang đọc những dòng này - đang khoái chí vì cái cảnh tôi phải nhún nhảy đi tìm chỗ giải quyết. Cậu biết tôi là một cô gái có khả năng cảm thấu tâm lý phức tạp, biết khuyên nhủ người khác, mà nỡ lòng nào lại bắt tôi cầm cái bể nước trong bụng đi dạo giữa hư vô mà chỉ chờ được xả ra ư!?

Thật không thể tha thứ!

\dots

Còn cái tên tác giả kia? Ờ, cái người đang dựng câu chuyện này từ nãy tới giờ?

Không biết xấu hổ. Cậu làm cho Suzu vẫn chưa đủ sao, mà lại còn lôi cả tôi vào cái trò dị hợm này!?

Tôi thấy cậu đang ngồi trong phòng máy lạnh, bịt tai nghe, bật nhạc nền ``Big Wave Beach'' từ Plants vs. Zombies 2, làm bộ làm tịch như thể: ``Ai chửi mắng thì ta giả điếc\~{}''

Tôi biết hết! Cậu nghĩ tôi không thấy sao? Tôi ở đây khổ sở từng giây một, còn cậu ngồi đó cười như xem hài kịch à!? Đồ-!

*ROẸT!*

\dots Khỉ thật, tôi bị cắt khỏi tầm tiếp xúc với độc giả rồi! Các người sẽ phải trả giá!

Bực mình quá, tôi nhấc chân bước nhanh hơn, cố tìm một chỗ nào đó khuất khuất một tí để giải quyết, thì bất chợt\dots

\il

``Ái\dots!''

Tôi vấp phải một cái gì đó dưới chân, loạng choạng suýt ngã.

``\vie{Ôi cái định công\dots!}'' tôi lảo đảo, may mắn thay đáp mông xuống đất.

``\vie{Sao đấy?}'' Kuon-kun quay lại, nhìn tôi đầy ngán ngẩm. ``\vie{Mắt mũi để ở chân à, Kaigou-chan?}''

Cái cách cậu ta chọc tôi cứ như thể tôi là đứa trẻ học đi vậy. Nếu bình thường thì tôi đã sừng cổ gào lên rồi, nhưng\dots

Tôi sững lại.

Mắt tôi nhìn về phía trước. Một cái gì đó không giống những thứ khác.

Một cánh cửa.

Đúng. Một cánh cửa gỗ cũ kỹ, đứng trơ trọi giữa vùng sương mù tím. Không có tường. Không có nhà. Không có bất kỳ cái gì khác xung quanh. Chỉ là một cánh cửa đứng trơ trọi.

Không tay nắm cửa Không biển hiệu. Không khóa. Không tiếng động.

Không gì cả.

Chúng tôi chưa từng thấy cái gì như vậy kể từ lúc bước chân vào nơi này.

Tôi nuốt nước bọt, cơn buồn đi vệ sinh tạm thời quên sạch.

``K-Kuon-kun\dots\ Suzu\dots'' tôi khẽ gọi.

Cậu ấy đã thấy nó, nét mặt thoáng ngạc nhiên.

Suzu cũng ngẩng đầu lên, đôi mắt khẽ trố ra.

Ba chúng tôi đứng đó, không ai nói gì. Chỉ còn tiếng tim tôi đập thình thịch, và một suy nghĩ lởn vởn trong đầu: Chuyện quái gì\dots\ đang chờ sau cánh cửa kia?

\pov{Hikawa Kuon}

Cánh cửa bằng gỗ cũ kỹ đứng sừng sững trước mặt chúng tôi như thể ai đó đã đặt nó ở đây từ lâu lắm rồi. Không có bất kỳ dòng chữ, ký hiệu hay dấu hiệu nào trên đó, chỉ là một cánh cửa đơn độc, trơ trọi giữa khoảng không vô tận của The Void.

Tôi bước tới, tay đặt lên cánh cửa, rồi nhìn Kaigou-chan và Suzuki.

``\vie{Chuẩn bị chưa?}'' tôi hỏi nhỏ.

Kaigou-chan gật đầu, vẻ mặt cảnh giác nhưng kiên định. Cô em gái nắm nhẹ lấy vai tôi, như để tự trấn an chính mình.

Tôi hít một hơi thật sâu, rồi đẩy nhẹ.

Một luồng ánh sáng xoáy cuộn trào ra từ phía trong, giống như một vòng xoáy dịch chuyển không ngừng: chuyển động tròn, uốn lượn như nước, với những vệt sáng đan xen vào nhau, tạo nên cảm giác vừa huyễn hoặc, vừa thôi miên. Tôi cảm thấy lạnh gáy ngay khi ánh mắt chạm vào nó. Một thứ gì đó rất giống thứ mà tôi và Kaigou từng thấy ở thế giới của mình.

Không\dots\ chính là nó. Thứ quái quỷ đã đưa chúng tôi rơi vào cái nơi chết tiệt này.

Tôi cúi người xuống, thử đưa chân ra phía trước, nhưng bên dưới chỉ là khoảng không. Không có mặt đất, không có điểm tựa. Chỉ cần thêm một chút thôi là tôi sẽ rơi hẳn vào đó.

Cô nàng đứng cạnh tôi thở dài: ``\vie{Không có lối đi khác, phải không?}''

Tôi lắc đầu. Tôi không nghĩ ở The Void có thể có đường nào khác nữa.

Suzuki nói nhẹ: ``\vie{Nhưng nếu cái này giống với cái ở thế giới cũ, thì có thể\dots\ nó là một dạng cổng dịch chuyển?}''

``\vie{Cổng giữa các chiều không gian?}'' tôi khẽ lặp lại lời giải thích của con bé. ``\vie{Hoặc chí ít\dots\ là một thứ giúp ta thoát khỏi The Void.}''

Kaigou-chan bồi thêm, dường như đồng tình với tôi: ``\vie{Cậu thấy đấy, nó không giống mấy cái bẫy. Nếu muốn hại ta thì nó đã làm rồi. Nhưng giờ thì nó chỉ đứng yên chờ ta bước vào.}''

Tôi gật đầu. ``\vie{Và nếu để ý kỹ\dots\ thì hình dáng của nó gần như trùng khớp với cái mà chúng ta từng thấy khi bị hút vào đây.}''

Chúng tôi im lặng mất vài giây. Ai cũng đang cố xâu chuỗi lại tất cả những gì đã xảy ra từ lúc bị hút vào, đến khi gặp nhau, rồi trò chuyện, đấu tranh, và cùng nhau bước đi. Hỉ nộ ái ố, tất cả mọi cung bậc cảm xúc, ba người chúng tôi đều đã trải qua đủ.

Và giờ đây, cánh cửa phía trước chúng tôi có lẽ sẽ đưa chúng tôi tới một nơi nào đó.

Có thể là một thế giới nào đó mà tôi không hề biết tới.

Có thể là thế giới của Kaigou-chan và Suzuki.

Hoặc\dots\ có thể là một nơi gọi là ``nhà.''

``\vie{Vậy thì\dots}'' tôi nói, quay sang hai người bạn đồng hành, ``\vie{\dots chúng ta sẽ cùng đi tiếp chứ? Kaigou-chan? Suzuki?}''

Người chị đưa ngón tay cái lên. Cô em gái mỉm cười, dẫu vẫn còn hơi run. Cả hai người họ đã sẵn sàng.

Tôi giơ tay lên. Không còn chờ gì nữa. Không còn thứ gì để chúng tôi có thể ở lại The Void này thêm một giây phút nào nữa.

``Một\dots''

``Hai\dots''

``\textbf{Ba!}''

Không ai ngập ngừng. Cả ba chúng tôi cùng siết tay nhau, rồi nhảy vào.

Cảm giác đầu tiên là gió. Gió lạnh và mạnh như một cơn bão đang lôi chúng tôi xuống đáy sâu hun hút của đại dương ánh sáng.

Không có tiếng la hét. Không bộc lộ nỗi sợ hãi, như lần đầu mà chúng tôi trải nghiệm. Chỉ là một cảm giác rất thật\dots\ rằng mình đang bị kéo vào một chiều không gian khác. Một lực hút vô hình mạnh khủng khiếp xoáy chúng tôi vào vùng tâm của luồng sáng đó. Cảm giác choáng váng, bồng bềnh, mọi giác quan đều bị đảo lộn.

Suzuki ôm chặt vào lưng tôi hơn, còn Kaigou-chan siết lấy tay tôi, như thể cả ba chúng tôi đang cố giữ lấy nhau khỏi tan biến trong luồng xoáy vô định ấy.

Không biết nơi đó sẽ đưa chúng tôi đi đâu. Không biết phía bên kia có gì đang chờ đợi.

Chỉ biết rằng\dots

Chúng tôi đã rời khỏi The Void.

Một hành trình mới của nhà Hikawa chúng tôi\dots\ vừa mới bắt đầu.

\end{document}